\chapter{Wireline and Optical Communications}
\label{ch:wireline}

\begin{chapterintro}
Radio gets the headlines, but almost every bit you send ends up on a wire or a fibre within a
few hundred metres of leaving your antenna. The cellular base station is fed by fibre, the
Wi-Fi access point by a twisted pair, and more than nine-tenths of intercontinental traffic
crosses the oceans on a few hundred submarine cables. This chapter is about those guided
channels. It starts with the transmission-line physics that every wireline engineer carries
around, then follows the telephone network from the hand-cranked magneto to the stored-program
switch, the signalling network that made it intelligent and the voiceband modems that pushed
it to its Shannon limit. It revisits the digital hierarchy from T1 through SONET to the optical
transport network, and then shows how DSL, cable modems and Ethernet squeezed gigabits out of
copper that was laid for telephone calls and analog television. The second half turns to
light: the physics of the fibre (loss, dispersion and nonlinearity), lasers, modulators and
photodiodes, the erbium-doped fibre amplifier and wavelength-division multiplexing, and the
coherent revolution that brought the QAM, MIMO equalisers and soft-decision codes of earlier
chapters into the optical domain. It closes with passive optical networks, data-centre optics
and a short look at free-space optics. Much of the book meets here: line codes
(Chapter~\ref{ch:baseband}), equalisation and echo cancellation (Chapter~\ref{ch:equalization}),
DMT (Chapter~\ref{ch:ofdm}), capacity (Chapter~\ref{ch:infotheory}) and modern codes
(Chapter~\ref{ch:moderncodes}) are all at work in the systems of this chapter.
\end{chapterintro}

%======================================================================
\section{Why guided media still carry the world}
\label{sec:ch24:intro}
%======================================================================

A wireless channel is shared, power-limited and unpredictable. A guided channel---a twisted
pair, a coaxial cable, an optical fibre---is private, stable and, in the case of fibre,
almost unimaginably wide. A single modern fibre pair carries tens of terabits per second
across an ocean; the entire radio spectrum below 100\,GHz, used perfectly, would carry less
than that over a single link. That is why the architecture of every network, wireless ones
included, is the same: radio for the last few hundred metres, where mobility matters, and
guided media everywhere else.

The guided world splits naturally into three tiers, and this chapter follows them.
\begin{description}[style=nextline,leftmargin=1.2em,font=\sffamily\bfseries\small\color{navy}]
\item[The core.] Long-haul and submarine fibre, carrying aggregated traffic between cities
and continents. The engineering problem is to maximise capacity times distance per fibre and
per watt. Coherent detection, optical amplifiers and wavelength-division multiplexing rule
here.
\item[The metro and data centre.] Fibre rings between central offices and links of a few
metres to a hundred kilometres between switches and servers. The problem is cost and power
per bit at enormous volume: hundreds of millions of optical transceivers a year.
\item[The access network, or ``last mile''.] The final connection to homes and businesses.
For a century it was a copper pair per subscriber; then coaxial cable built for television;
now increasingly fibre. The problem is cost per subscriber, and the defining constraint has
been reuse: the copper was already in the ground, and anything that could squeeze more out of
it without digging was worth enormous engineering effort.
\end{description}

\mdcfig[0.85]{ch24_access_rates}{Approximate maximum downstream rates of access technologies
against the year their standard (or first major version) appeared. Each family advanced by
reusing the same physical medium with more signal processing; together they track a growth of
roughly 50\% per year, an observation sometimes called Nielsen's law. Years and rates are
approximate and are for orientation only.}

Figure~\ref{fig:ch24_access_rates} shows the result of that pressure. Over sixty years the
rate delivered to a home grew by a factor of about $10^8$, from 300\,b/s acoustic-coupled
modems to 10\,Gb/s passive optical networks. Remarkably, much of that growth came over media
that did not change: the telephone pair that carried a 300\,b/s modem in 1962 carried
100\,Mb/s of VDSL2 in 2010. What changed was signal processing---equalisation, echo
cancellation, multicarrier modulation, coding and, finally, crosstalk cancellation. The
copper story is therefore a long case study in the ideas of this book, and the fibre story
is the same ideas applied to a channel a million times wider.

%======================================================================
\section{Transmission lines for communication engineers}
\label{sec:ch24:tline}
%======================================================================

Every wire is a transmission line once its length is comparable to a wavelength, or, more
practically, once the signal's rise time is comparable to the propagation delay. A telephone
pair is a transmission line at a few kilohertz over several kilometres; a circuit-board trace
is one at a few gigahertz over a few centimetres. This section collects the results a
communication engineer needs; the electromagnetic derivations are in any microwave text.

\subsection{The RLGC model and the telegrapher's equations}

Model a uniform two-conductor line as a cascade of infinitesimal sections, each of length
$dz$, with series resistance $R\,dz$ and inductance $L\,dz$, and shunt conductance $G\,dz$ and
capacitance $C\,dz$. $R$, $L$, $G$ and $C$ are the \term{primary line constants}, per unit
length. Kirchhoff's laws applied to one section give, for phasors at angular frequency
$\omega$, the \term{telegrapher's equations}
\begin{equation}
\frac{dV}{dz}=-(R+j\omega L)\,I,\qquad \frac{dI}{dz}=-(G+j\omega C)\,V,
\label{eq:ch24:telegrapher}
\end{equation}
first written down in this form by Oliver Heaviside in the 1880s. Eliminating $I$ gives the
wave equation $d^2V/dz^2=\gamma^2V$, whose solution is a pair of waves travelling in
opposite directions,
\begin{equation}
V(z)=V^+e^{-\gamma z}+V^-e^{+\gamma z},\qquad
I(z)=\frac{1}{Z_0}\left(V^+e^{-\gamma z}-V^-e^{+\gamma z}\right),
\end{equation}
with the \term{propagation constant} and the \term{characteristic impedance}
\begin{equation}
\gamma=\alpha+j\beta=\sqrt{(R+j\omega L)(G+j\omega C)},\qquad
Z_0=\sqrt{\frac{R+j\omega L}{G+j\omega C}}.
\label{eq:ch24:gammaz0}
\end{equation}
The real part $\alpha$ is the attenuation in nepers per metre (multiply by 8.686 for decibels);
the imaginary part $\beta$ is the phase constant, and the phase velocity is $\omega/\beta$. A
line of length $\ell$ terminated in its own characteristic impedance looks, from its input,
like a resistor $Z_0$ and transfers the signal as $H(f)=e^{-\gamma\ell}$: an attenuation of
$8.686\,\alpha\ell$\,dB and a delay of $\beta\ell/\omega$.

Two limits explain almost everything about copper.

\emph{At low frequencies}, $\omega L\ll R$ and $G\approx0$, so $\gamma\approx\sqrt{j\omega RC}$
and $Z_0\approx\sqrt{R/(j\omega C)}$. The attenuation grows as $\sqrt f$, the impedance is large
and has a phase of $-45^\circ$, and the line behaves as a diffusion medium. This is Kelvin's
cable of Chapter~\ref{ch:history}, where the response time grows as the square of the length.
A telephone pair is in this regime throughout the voice band: a 26-gauge pair has
$|Z_0|\approx970$\,$\Omega$ at 1\,kHz (Figure~\ref{fig:ch24_line_loss}), which is why
telephone circuits were designed around 600 or 900\,$\Omega$ terminations.

\emph{At high frequencies}, $\omega L\gg R$ and $\omega C\gg G$, and a first-order expansion gives
\begin{equation}
Z_0\approx\sqrt{\frac LC},\qquad
\alpha\approx\frac{R}{2Z_0}+\frac{GZ_0}{2},\qquad v=\frac{1}{\sqrt{LC}}.
\label{eq:ch24:hfloss}
\end{equation}
The impedance becomes real and nearly constant---about 100\,$\Omega$ for twisted pair,
50 or 75\,$\Omega$ for coax---and the line becomes a proper waveguide with a velocity of 0.6
to 0.8 times that of light. The first term of $\alpha$ is conductor loss and the second is
dielectric loss.

\subsection{Skin effect, dielectric loss and the shape of copper loss}

At high frequency the current crowds into a skin of depth $\delta=1/\sqrt{\pi f\mu\sigma}$ at
the conductor surface (66\,$\mu$m in copper at 1\,MHz, 6.6\,$\mu$m at 100\,MHz), so $R$ rises as
$\sqrt f$ and the conductor term of~\eqref{eq:ch24:hfloss} gives a loss in decibels
proportional to $\sqrt f$. The dielectric term $GZ_0/2$, with $G=\omega C\tan\delta$ for an
insulator of loss tangent $\tan\delta$, grows in proportion to $f$. The loss of any copper
medium is therefore well fitted by
\begin{equation}
\alpha_{\mathrm{dB}}(f)\approx k_1\sqrt f+k_2f,
\label{eq:ch24:k1k2}
\end{equation}
with skin effect dominating in telephone cable and coax, and dielectric loss taking over on
circuit boards at gigahertz frequencies (Chapter~\ref{ch:baseband}). Cable data sheets and the
TIA-568 cabling standards specify insertion loss in exactly this form, sometimes with a third
$k_3/\sqrt f$ term for low-frequency effects.

\mdcfig{ch24_line_loss}{Left: attenuation per kilometre of 26 and 24 AWG telephone pairs
(from the parametric RLGC model used in DSL standards work), a Category~6 data pair (from
the TIA-568 insertion-loss formula) and typical RG-6 drop coax. In the voice band the loss
grows as $\sqrt f$ through the RC diffusion mechanism; above about 100\,kHz the skin effect
takes over, again as $\sqrt f$; dielectric loss adds a component proportional to $f$. Right:
the magnitude (solid) and phase (dashed) of the characteristic impedance of a telephone pair,
falling from the RC regime (about 700--1000\,$\Omega$ at $-45^\circ$) to a nearly real
100\,$\Omega$ above about 100\,kHz.}

Figure~\ref{fig:ch24_line_loss} shows the result for the four copper media of this chapter.
Note the logarithmic axes: at 1\,MHz a 26-gauge telephone pair loses about 25\,dB per
kilometre, and at 30\,MHz about 150\,dB/km. No conceivable signal processing extracts data
from a channel with hundreds of decibels of loss, which is why every increase in copper
access speed has been bought by shortening the copper.

\begin{worked}[The loss of a 3\,km loop at 1\,MHz]
A subscriber is served by 3\,km of 26 AWG (0.4\,mm) cable. Estimate the loss at 1\,MHz.

At 1\,MHz the parametric cable model gives, per kilometre, $R\approx627$\,$\Omega$,
$L\approx0.57$\,mH, $C\approx49$\,nF and a negligible $G$. Since $\omega L=3.6$\,k$\Omega$ per
km is six times $R$, the high-frequency approximation applies: $Z_0\approx\sqrt{L/C}\approx
108$\,$\Omega$, and
\[
\alpha\approx\frac{R}{2Z_0}=\frac{627}{2\times108}\approx2.9\ \text{Np/km}
\;\Rightarrow\;8.686\times2.9\approx25\ \text{dB/km}.
\]
The full expression~\eqref{eq:ch24:gammaz0} gives 25.4\,dB/km. Three kilometres therefore lose
about 76\,dB at 1\,MHz. If the transmitter launches $-40$\,dBm/Hz (the ADSL downstream level)
and the receiver sees the conventional background noise of $-140$\,dBm/Hz, the SNR at 1\,MHz is
$100-76=24$\,dB before any crosstalk---enough for about 5 bits per tone after the SNR gap of
Chapter~\ref{ch:ofdm}, and falling by about 1\,bit for every 6\,dB of extra loss. The velocity
is $v=1/\sqrt{LC}\approx1.9\times10^8$\,m/s, about $0.63c$, so the loop delay is 16\,$\mu$s.
\end{worked}

\subsection{Twisted pair and coaxial cable}

A \term{twisted pair} is two insulated conductors twisted together. The twist is the whole
point: an interfering magnetic field induces voltages of opposite sign in successive half-twists,
which cancel, and the two wires of the pair are, on average, equally exposed to any nearby
conductor, so a \emph{differential} receiver rejects what they pick up in common. The cancellation
is never perfect, and the residue is crosstalk. Telephone cables bundle pairs into
\emph{binder groups} of 25 pairs, each pair with a different twist length so that neighbours do
not lie parallel for long, and bundle binders into cables of up to several thousand pairs.
Conductor sizes are given in the American Wire Gauge: 26 AWG is 0.40\,mm, 24 AWG 0.51\,mm, 22 AWG
0.64\,mm, and 19 AWG 0.91\,mm. Thinner wire is cheaper and was used close to the exchange where
cables are biggest; thicker wire was used on long rural loops.

Data cabling follows the TIA-568 and ISO/IEC 11801 categories: Category~3 (16\,MHz, 10BASE-T),
Category~5 and 5e (100\,MHz, Fast and Gigabit Ethernet), Category~6 (250\,MHz), Category~6A
(500\,MHz, 10GBASE-T to 100\,m) and Category~8 (2\,GHz, 25/40GBASE-T to 30\,m). Each category
tightens the twist and the manufacturing tolerances, and the higher ones add foil shields.

A \term{coaxial cable} confines the field between an inner conductor and a surrounding shield,
so it neither radiates nor picks up interference as long as the shield is intact, and its loss
is set almost entirely by the skin effect in the conductors. Cable television uses 75\,$\Omega$
coax, a value that minimises loss for an air or foam dielectric; laboratory and radio work uses
50\,$\Omega$, a compromise between minimum loss and maximum power handling. A 1\,cm aluminium
``hardline'' trunk cable loses about 2\,dB per 100\,m at 500\,MHz; an RG-6 drop cable to the house
loses about 15\,dB per 100\,m at the same frequency.

\subsection{Reflections and return loss}

When a line of impedance $Z_0$ is terminated in a load $Z_L\ne Z_0$, part of the incident wave
is reflected. The \term{reflection coefficient} and the \term{return loss} are
\begin{equation}
\Gamma=\frac{Z_L-Z_0}{Z_L+Z_0},\qquad \mathrm{RL}=-20\log_{10}|\Gamma|\ \text{dB}.
\label{eq:ch24:gamma}
\end{equation}
A 100\,$\Omega$ line terminated in 120\,$\Omega$ has $\Gamma=0.09$ and a return loss of 21\,dB; a
75\,$\Omega$ coax connected to a 50\,$\Omega$ instrument has $\Gamma=-0.2$ and 14\,dB. An open
circuit reflects everything with $\Gamma=+1$, a short with $\Gamma=-1$. The \term{time-domain
reflectometer} (TDR), which launches a fast step and watches the reflections come back, turns
this into a tool: the delay of each echo locates a fault, and its sign and size identify it. Every
line tester in a telephone company's van is a TDR.

Reflections matter to a communication system in two ways. First, they are \emph{echoes}: in a
full-duplex system the transmitter's own signal, reflected from the near-end impedance mismatch,
arrives at the local receiver, often much stronger than the far-end signal (Section~\ref{sec:ch24:hybrid}).
Second, a mismatch partway along a line---a change of gauge, a connector, a branch---produces a
multipath channel. The most notorious example in the telephone plant is the \term{bridged tap}:
an unused branch of cable left connected to a loop, a legacy of the practice of wiring one pair to
several possible customer locations and connecting whichever was needed. An open-ended tap of
length $\ell_t$ presents at its junction the admittance $Y=\tanh(\gamma\ell_t)/Z_0$; at the
frequencies where the tap is an odd number of quarter wavelengths long, $f=(2k+1)v/(4\ell_t)$, the
reflected wave returns in antiphase and the tap shorts the line (Figure~\ref{fig:ch24_loop_impairments}).
A 25\,m tap carves its first notch near 2\,MHz, squarely in the VDSL2 band.

\subsection{Crosstalk}

Crosstalk is the coupling of signals between pairs in the same cable, through the mutual
capacitance and inductance of imperfectly twisted neighbours. It comes in two kinds, named for
where the victim's receiver is relative to the disturbing transmitter.
\begin{description}[style=nextline,leftmargin=1.2em,font=\sffamily\bfseries\small\color{navy}]
\item[Near-end crosstalk (NEXT).] The disturber's transmitter and the victim's receiver are at
the same end. The coupled signal has not travelled down the cable, so it is strong, and it
grows with frequency, because coupling through a small capacitance rises as $f$ and the
contributions along the line add with random phase. The standard empirical model, due to
Unger and adopted in the ANSI and ITU DSL standards, gives the power coupling from $N$
disturbers in a 50-pair binder as
\begin{equation}
|H_{\mathrm{NEXT}}(f)|^2=K_{\mathrm{NEXT}}\left(\frac{N}{49}\right)^{0.6}f^{3/2},\qquad
K_{\mathrm{NEXT}}\approx8.8\times10^{-14},
\label{eq:ch24:next}
\end{equation}
with $f$ in hertz. It is a ``1\% worst case'' model: 99\% of pairs in real cables do better. At
80\,kHz, the ISDN frequency, it gives 57\,dB of NEXT loss for 49 disturbers, the number on which
ISDN was designed. The $N^{0.6}$ law reflects the fact that the nearest neighbours dominate.
\item[Far-end crosstalk (FEXT).] The disturber's transmitter is at the far end, with the
victim's. The coupling accumulates along the whole length and travels with the signal, so it is
attenuated by the line like the signal itself:
\begin{equation}
|H_{\mathrm{FEXT}}(f)|^2=K_{\mathrm{FEXT}}\left(\frac{N}{49}\right)^{0.6}f^{2}\,\ell\,|H(f)|^2,
\qquad K_{\mathrm{FEXT}}\approx7.7\times10^{-21}\ (\ell\text{ in feet}).
\label{eq:ch24:fext}
\end{equation}
\end{description}
Figure~\ref{fig:ch24_crosstalk_vectoring} plots both against the through signal of a 1\,km loop. NEXT
exceeds the received signal above about 2\,MHz, so systems that transmit in both directions in the
same band on adjacent pairs cannot work there. That single observation dictates the architecture
of DSL: ADSL and VDSL2 use \emph{frequency-division duplexing} with different bands for each
direction, so that NEXT from other lines falls outside the receiver's band, and G.fast uses
synchronised time-division duplexing, so that no receiver listens while a neighbour transmits.
What remains is FEXT, which, unlike the signal, \emph{grows} relative to it as $f^2\ell$, and which
limits every modern DSL system until it is cancelled (Section~\ref{sec:ch24:vectoring}).

\begin{pitfall}[Self-crosstalk is the real enemy]
Engineers new to DSL look first at background noise, typically $-140$\,dBm/Hz, and compute
spectacular rates. On a real cable the dominant noise is crosstalk from other lines carrying
\emph{the same service}---self-NEXT and self-FEXT---because the disturbers use exactly the victim's
band at exactly the victim's power. Adding subscribers to a binder therefore lowers everyone's
rate, and a rate promised from a single-line test on an empty cable will not be achieved in
service. Always include a crosstalk model with a realistic number of same-service disturbers;
the Unger model with 10 to 24 disturbers is the usual planning choice.
\end{pitfall}

%======================================================================
\section{The telephone plant}
\label{sec:ch24:telephone}
%======================================================================

\subsection{From magneto to stored program control}

\begin{history}[A hundred years of switching]
Bell's telephone (1876) needed a way to connect any subscriber to any other, and the first
answer was the manual switchboard, opened in New Haven, Connecticut, in 1878 with 21 subscribers.
Operators plugged cords into jacks; the subscriber's telephone carried its own battery
and a hand-cranked magneto to ring the exchange. By the 1890s \emph{common battery} working put the
battery at the exchange, where it still is: the $-48$\,V on a telephone line today descends from
those lead-acid cells.

Automation began with Almon Strowger's step-by-step switch, patented in 1891 and first used at La
Porte, Indiana, in 1892. Each digit dialled moved an electromechanical selector one step per pulse;
the call's path through the exchange was built digit by digit, directly controlled by the
subscriber. The Bell System's \emph{panel} switches of the 1920s and the \emph{crossbar}
switches (the No.~1 crossbar entered service in 1938) separated control from switching: common
control equipment received all the digits first, then chose and set up a path through a matrix of
relay contacts. Crossbar offices carried most of North America's calls into the 1970s.

The decisive step was Bell Labs' No.~1 Electronic Switching System, cut over in Succasunna, New
Jersey, in 1965: the switching fabric was still made of metallic reed relays, but the control was a
\emph{stored-program} computer. Features became software. Call waiting, call forwarding and
three-way calling arrived as program changes, and stored program control has defined telephone
switching ever since. The No.~4 ESS (1976) was the first digital toll switch, switching PCM time
slots instead of metallic paths; the No.~5 ESS (1982) and Northern Telecom's DMS-100 brought
digital switching to local exchanges. A digital switch connects calls by copying 8-bit samples
between time slots in memory, a technique called \emph{time-slot interchange}.
\end{history}

\subsection{The local loop}

The \term{local loop} is the pair of copper wires between the exchange (the \emph{central office},
CO) and the customer. It is the largest single investment of any telephone company: several hundred
million pairs worldwide, each running on average a few kilometres. Leaving the CO, pairs travel in
large \emph{feeder} cables to a cross-connect cabinet (in North America the \emph{serving area
interface}), then in smaller \emph{distribution} cables along the streets, and finally in a
\emph{drop} wire to the house. North American loops are typically 1 to 6\,km long, with a long tail
of rural loops beyond 8\,km; European and Asian loops are generally shorter.

The CO powers the line with about $-48$\,V through a feed resistance; when the handset goes off-hook,
a loop current of 20 to 80\,mA flows, which the exchange detects. Dialling interrupts the current
(pulse dialling at 10 pulses per second) or sends dual-tone multi-frequency (DTMF) pairs; ringing
is a high-voltage AC signal, about 90\,V\,rms at 20\,Hz in North America, superimposed on the battery.
The \term{line card} at the exchange implements all of this. Its functions are summarised by the
mnemonic BORSCHT: \emph{B}attery feed, \emph{O}vervoltage protection, \emph{R}inging,
\emph{S}upervision (on/off-hook detection), \emph{C}odec (the $\mu$-law or A-law PCM converter of
Chapter~\ref{ch:pcm}), \emph{H}ybrid (two-wire to four-wire conversion) and \emph{T}esting.

\begin{figure}[htbp]\centering
\begin{tikzpicture}[
  blk/.style={draw=navy,fill=navy!6,thick,minimum height=0.8cm,minimum width=1.3cm,align=center,font=\sffamily\scriptsize},
  arr/.style={-{Stealth[length=2mm]},thick,navy}]
\node[blk,minimum height=2.4cm] (phone) at (0,0) {Tele-\\phone};
\node[blk,minimum height=2.4cm,fill=accent!6,draw=accent] (hyb) at (3.6,0) {Hybrid\\[2pt]balance\\network\\$Z_B$};
\draw[thick,navy] (0.65,0.2) -- (2.95,0.2); \draw[thick,navy] (0.65,-0.2) -- (2.95,-0.2);
\node[font=\sffamily\scriptsize,text=navy,align=center] at (1.8,0.65) {2-wire loop};
\node[font=\sffamily\scriptsize,text=navy,align=center] at (1.8,-0.65) {both directions};
\node[blk] (adc) at (6.0,0.85) {A/D\\(codec)};
\node[blk] (dac) at (6.0,-0.85) {D/A\\(codec)};
\node[draw=navy,thick,circle,inner sep=1pt,font=\small] (sum) at (8.0,0.85) {$+$};
\node[blk] (ec) at (8.0,-0.0) {Echo\\canceller};
\node[blk,minimum width=1.7cm] (sw) at (10.4,0.85) {to network\\(64\,kb/s)};
\node[blk,minimum width=1.7cm] (sw2) at (10.4,-0.85) {from network\\(64\,kb/s)};
\draw[arr] (hyb.east |- adc) -- (adc); \draw[arr] (adc) -- (sum); \draw[arr] (sum) -- (sw);
\draw[arr] (sw2) -- (dac); \draw[arr] (dac) -- (hyb.east |- dac);
\draw[arr] (9.1,-0.85) |- (ec.east);
\draw[arr] (ec) -- node[right,font=\scriptsize]{$-$} (sum);
\draw[arr,dashed,accent] (4.2,-0.75) .. controls (4.75,-0.3) and (4.75,0.3) .. node[right,font=\sffamily\scriptsize,text=accent]{echo} (4.2,0.75);
\node[draw=navy!50,dashed,rounded corners,fit=(hyb)(adc)(dac)(ec),inner sep=0.2cm,label={[font=\sffamily\scriptsize,text=navy]above:line card (four-wire side to the right of the hybrid)}] {};
\end{tikzpicture}
\caption{Two-wire to four-wire conversion. The loop carries both directions on one pair; the hybrid
separates them using a balance impedance $Z_B$ that mimics the line. Because $Z_B$ never matches the
real loop exactly, part of the incoming signal from the network leaks back towards it as echo; an
adaptive echo canceller (Chapter~\ref{ch:equalization}) removes the residue.}
\label{fig:ch24_hybrid}
\end{figure}

\subsection{The hybrid, two wires and four}
\label{sec:ch24:hybrid}

A telephone loop carries both directions of speech on one pair, but amplifiers, multiplexers and
digital transmission are one-way devices. Somewhere, therefore, every long-distance circuit must
convert from \emph{two-wire} to \emph{four-wire} working, with a separate path for each direction.
The device that does it is the \term{hybrid} (Figure~\ref{fig:ch24_hybrid}), originally a
transformer bridge and now an operational-amplifier circuit inside the line card. It is a balanced
bridge: the signal from the network drives the line and a \emph{balance impedance} $Z_B$ in
opposite arms, and the receive port sees the difference. If $Z_B$ equals the loop's input impedance,
the transmitted signal cancels at the receive port exactly, while the signal arriving from the
loop does not. The figure of merit is the \term{transhybrid loss}, the attenuation from the
four-wire input to the four-wire output, which is related to the mismatch by the same formula as
return loss,
\begin{equation}
\text{THL}\approx-20\log_{10}\left|\frac{Z_{\text{in}}-Z_B}{Z_{\text{in}}+Z_B}\right|.
\end{equation}
The trouble is that $Z_{\text{in}}$ depends on the gauge, length, bridged taps and loading of each
individual loop (Figure~\ref{fig:ch24_line_loss} shows how strongly $Z_0$ varies in the voice
band), while $Z_B$ is one compromise network for all loops. A typical hybrid achieves only 6 to
20\,dB of transhybrid loss, with an average around 11\,dB.

The leakage is \term{echo}. On a local call the talker's own voice returns within a few
milliseconds and sounds like the natural sidetone of the handset. As the round-trip delay grows,
echo becomes annoying and then intolerable: ITU-T G.131 relates the acceptable echo level to the
delay, and above roughly 25\,ms of one-way delay control is required. Early long-distance circuits
used \emph{echo suppressors}, voice-activated switches that opened the return path while the far
party talked, at the cost of clipped speech whenever both spoke. Satellite circuits, with 270\,ms
one-way delay (Chapter~\ref{ch:satellite}), forced the invention of the adaptive \term{echo
canceller}, developed at Bell Labs in the 1960s by Kelly, Logan and Sondhi; the first single-chip
version followed in 1980. Modern network echo cancellers (ITU-T G.168) model echo paths up to
64 or 128\,ms long with NLMS adaptive filters (Section~\ref{sec:ch12:echo}). A second hazard of the
hybrid is \emph{singing}: if the loop gain around the four-wire circuit---two amplifiers and two
transhybrid losses---exceeds unity, the circuit oscillates, which is why analog four-wire circuits
always included some loss.

\begin{keyidea}[Echo is the price of sharing a wire]
Whenever two directions share one medium at the same frequencies, each receiver hears its own
transmitter. The hybrid, the echo canceller of V.32, the four echo-and-NEXT cancellers of
Gigabit Ethernet, DOCSIS~4.0's full-duplex cable plant and the self-interference cancellation of
full-duplex radio are the same idea at successively higher bandwidths: model the leakage from your
own known transmit signal and subtract it. The alternative is to divide the medium in frequency or
time, which wastes half of it.
\end{keyidea}

\subsection{Loading coils}

The voice-band loss of a long unloaded loop rises with frequency (the RC regime of the previous
section), which muffles speech. Heaviside showed that adding inductance moves the line towards his
distortionless condition $R/L=G/C$, and from about 1900 the Bell System inserted \term{loading
coils} into long loops. The North American standard H88 loading places 88\,mH coils every 6000\,ft
(1.83\,km), with the first coil 3000\,ft from the exchange; loops longer than about 18\,000\,ft
(5.5\,km) were loaded. The coils turn the cable into a lumped low-pass $LC$ ladder with a cutoff near
$f_c=1/(\pi\sqrt{L_{\mathrm{coil}}C_{\mathrm{section}}})$, which for H88 on 26-gauge cable is about
3.6\,kHz. Below the cutoff the loss is flat and lower than the unloaded loss; above it the loss
rises extremely steeply (Figure~\ref{fig:ch24_loop_impairments}).

\mdcfig{ch24_loop_impairments}{Two legacies of the telephone plant, computed with the two-port
(ABCD) loop model. Left: 5.5\,km of 26 AWG cable between 600\,$\Omega$ terminations, unloaded and
with H88 loading coils (88\,mH every 1.83\,km). Loading flattens the voice band and then cuts off
sharply near 3.6\,kHz, which made long loops good for speech and useless for DSL. Right: a
300\,m, 24 AWG loop with and without a bridged tap. Each open-ended tap produces notches at the
frequencies where it is an odd number of quarter wavelengths long; the shorter the tap, the higher
the first notch.}

Loading was a triumph for telephony and a disaster for data. A loaded loop passes nothing above
4\,kHz, so before any DSL service could be offered the coils had to be found and removed, a job of
digging up manholes that still appears in the cost estimates of rural broadband projects. Loading
also contributed the familiar 3.4\,kHz upper edge of the telephone channel and therefore, through
the 8\,kHz sampling rate of PCM, much of the timing of the digital network (Chapter~\ref{ch:pcm}).
The North American \emph{carrier serving area} rules of the 1980s, which define loops suitable for
digital services, specify no loading coils, at most about 12\,000\,ft of 24-gauge or 9000\,ft of
26-gauge cable, and limited bridged taps.

\subsection{The switched network and its hierarchy}

\begin{figure}[htbp]\centering
\begin{tikzpicture}[
  sw/.style={draw=navy,fill=navy!6,thick,rounded corners=2pt,minimum height=0.6cm,minimum width=1.5cm,align=center,font=\sffamily\scriptsize},
  stp/.style={draw=accent,fill=accent!6,thick,diamond,aspect=1.6,inner sep=1pt,font=\sffamily\scriptsize},
  ph/.style={draw=navy!60,fill=white,circle,inner sep=1.3pt,font=\sffamily\tiny}]
\node[sw] (c1a) at (0,3.6) {Class 1\\regional};
\node[sw] (c1b) at (5.2,3.6) {Class 1\\regional};
\node[sw] (c2a) at (0,2.6) {Class 2\\sectional};
\node[sw] (c2b) at (5.2,2.6) {Class 2\\sectional};
\node[sw] (c3a) at (0,1.6) {Class 3\\primary};
\node[sw] (c3b) at (5.2,1.6) {Class 3\\primary};
\node[sw] (c4a) at (0,0.6) {Class 4\\toll};
\node[sw] (c4b) at (5.2,0.6) {Class 4\\toll};
\node[sw] (c5a) at (-1.0,-0.6) {Class 5\\end office};
\node[sw] (c5b) at (1.0,-0.6) {Class 5\\end office};
\node[sw] (c5c) at (5.2,-0.6) {Class 5\\end office};
\foreach \a/\b in {c1a/c2a,c2a/c3a,c3a/c4a,c4a/c5a,c4a/c5b,c1b/c2b,c2b/c3b,c3b/c4b,c4b/c5c}
  \draw[thick,navy] (\a) -- (\b);
\draw[thick,navy] (c1a) -- node[above,font=\sffamily\scriptsize]{final route} (c1b);
\draw[thick,navy,dashed] (c4a) -- node[below,font=\sffamily\scriptsize,pos=0.5]{high-usage trunk} (c4b);
\draw[thick,navy,dashed] (c3a) -- (c4b);
\foreach \x in {-1.5,-1.0,-0.5} \draw[navy!60] (\x+0.0,-1.0) -- (\x,-1.4) node[ph,below]{};
\foreach \x in {0.5,1.0,1.5} \draw[navy!60] (\x,-1.0) -- (\x,-1.4) node[ph,below]{};
\foreach \x in {4.7,5.2,5.7} \draw[navy!60] (\x,-1.0) -- (\x,-1.4) node[ph,below]{};
\node[font=\sffamily\scriptsize,text=navy] at (2.6,-1.95) {subscriber loops};
% SS7 plane
\node[stp] (s1) at (9.0,2.7) {STP};
\node[stp] (s2) at (11.0,2.7) {STP};
\node[sw,fill=accent!6,draw=accent] (scp) at (10.0,3.8) {SCP\\(database)};
\node[sw] (ssp1) at (8.6,0.9) {SSP\\(switch)};
\node[sw] (ssp2) at (11.4,0.9) {SSP\\(switch)};
\draw[thick,accent] (s1) -- (s2); \draw[thick,accent] (s1) -- (scp); \draw[thick,accent] (s2) -- (scp);
\draw[thick,accent] (ssp1) -- (s1); \draw[thick,accent] (ssp1) -- (s2);
\draw[thick,accent] (ssp2) -- (s1); \draw[thick,accent] (ssp2) -- (s2);
\draw[thick,navy] (ssp1) -- (ssp2); \node[font=\sffamily\scriptsize,text=navy] at (10.0,0.25) {voice trunks};
\node[font=\sffamily\scriptsize\bfseries,text=navy] at (2.6,4.35) {voice network (Bell System, c.\ 1950--1984)};
\node[font=\sffamily\scriptsize\bfseries,text=accent] at (10.0,4.65) {SS7 signalling network};
\end{tikzpicture}
\caption{Left: the five-level switching hierarchy of the North American network before divestiture.
A call climbed only as far as needed to find a common switch, with high-usage trunks (dashed)
providing shortcuts that overflowed to the final route. Right: the out-of-band Signalling System
No.~7 network. Switches (service switching points) exchange signalling messages through duplicated
signal transfer points, and query service control points for number translation, mobility and
other services.}
\label{fig:ch24_pstn}
\end{figure}

The public switched telephone network (PSTN) is a hierarchy of switches connected by \emph{trunks}.
In the Bell System's plan, which lasted from the 1950s until the divestiture of 1984, there were five
classes of switching office (Figure~\ref{fig:ch24_pstn}): many thousands of class~5 end offices,
where the loops terminate, and above them toll, primary and sectional centres and, at the top, about
a dozen regional centres in the United States and Canada. A call climbed the
tree only as far as needed to reach a switch common to both parties, but heavily used pairs of offices
were also linked directly by \emph{high-usage trunks}, sized by Erlang's traffic formulas
(Chapter~\ref{ch:multipleaccess}), with calls overflowing to the hierarchical \emph{final route}
when the direct trunks were busy. From 1984 AT\&T replaced the fixed hierarchy in its long-distance
network with \emph{dynamic non-hierarchical routing}, in which the computer-controlled toll switches
chose among alternative two-hop paths according to the time of day and the measured load---an early
example of a network adapting its routing to traffic.

\subsection{Signalling: from 2600\,Hz to SS7}

A switch must tell the next switch who is being called and when a call starts and ends. For decades
this \term{signalling} travelled in the voice channel itself. Trunks used a 2600\,Hz tone to mark an
idle circuit and multi-frequency (MF) tone pairs to send digits. In-band signalling was cheap but had
a famous weakness: anything that could play 2600\,Hz into the mouthpiece could seize a trunk. The
``phone phreaks'' of the 1960s and 1970s did exactly that, first with a toy whistle and then with
electronic \emph{blue boxes}, and the episode is part of computing folklore.

The cure was to move signalling out of band onto a separate data network. AT\&T's Common Channel
Interoffice Signalling (CCIS, 1976) was followed by the CCITT's \term{Signalling System No.~7}
(SS7), specified in the ITU-T Q.700 series and deployed worldwide through the 1980s and 1990s.
SS7 is a packet network with its own protocol stack: the Message Transfer Part provides reliable,
duplicated delivery; the ISDN User Part (ISUP) sets up and releases calls; and the Transaction
Capabilities Application Part (TCAP) supports database queries. The nodes are \emph{service switching
points} (the telephone switches), duplicated \emph{signal transfer points} that route messages, and
\emph{service control points}, databases that translate toll-free numbers, support number
portability and, through the Mobile Application Part (MAP), keep track of every mobile phone's
location. SS7 made the telephone network programmable: the ``intelligent network'' of free-phone
numbers, calling cards and mobile roaming is built on it. It was designed for a closed club of
trusted operators, and its lack of authentication has, since the 2010s, allowed location tracking
and interception of text messages by anyone who obtains access; operators now filter SS7 traffic at
their borders, and 4G and 5G cores use the Diameter and HTTP-based protocols that replace it.

\begin{inpractice}[The PSTN today]
The circuit-switched network is being retired. Most operators have moved their trunks to voice over
IP, and many have announced dates in the late 2020s for switching off their copper PSTN altogether,
replacing analog lines with voice over broadband. The 64\,kb/s G.711 channel survives as the
universal fallback codec of VoIP (Chapter~\ref{ch:sourcecoding}); the 125\,$\mu$s timing of the
digital hierarchy survives in the synchronisation networks that time mobile base stations; and SS7
survives in the interconnection of mobile networks for roaming and text messages. The copper loops
themselves are being reused for DSL until fibre arrives, and in some places for powering
fibre-fed distribution points (Section~\ref{sec:ch24:gfast}).
\end{inpractice}

%======================================================================
\section{Voiceband modems}
\label{sec:ch24:modems}
%======================================================================

From the early 1960s to the end of the 1990s, the telephone network was also the world's data
network, and the voiceband modem was the device that adapted digital data to a channel designed
for speech: roughly 300 to 3400\,Hz, an SNR of 30 to 40\,dB on a good connection, and a
frequency-dependent delay and attenuation that varied from call to call. Modem history, outlined
in Chapter~\ref{ch:modulation}, is a compact tour of this book, and Table~\ref{tab:ch24:modems}
summarises it. Here we look at it from the wire's point of view.

\begin{table}[htbp]\centering\small
\caption{Landmark voiceband modem standards (rates are the maximum for each standard; years are
approximate dates of introduction or approval).}
\label{tab:ch24:modems}
\begin{tabularx}{\linewidth}{@{}l l l l X@{}}
\toprule
Standard & Year & Rate & Symbol rate & Technique\\
\midrule
Bell 103 / V.21 & 1962 & 300\,b/s & 300\,Bd & Binary FSK, frequency-split full duplex\\
Bell 212A / V.22 & 1980 & 1200\,b/s & 600\,Bd & Differential QPSK, frequency split\\
V.22bis & 1984 & 2400\,b/s & 600\,Bd & 16-QAM, frequency split, adaptive equaliser\\
V.29 & 1976 & 9600\,b/s & 2400\,Bd & 16-point constellation, four-wire leased lines and fax\\
V.32 & 1984 & 9600\,b/s & 2400\,Bd & 32-point cross with 8-state trellis code; echo cancellation\\
V.32bis & 1991 & 14.4\,kb/s & 2400\,Bd & 128-point TCM, echo cancellation\\
V.34 & 1994/96 & 33.6\,kb/s & up to 3429\,Bd & Line probing, precoding, shell mapping, 4D TCM\\
V.90 & 1998 & 56\,kb/s down & 8000/s & PCM levels downstream (digital server end)\\
V.92 & 2000 & 48\,kb/s up & 8000/s & PCM upstream, quick connect, modem on hold\\
\bottomrule
\end{tabularx}
\end{table}

\subsection{Full duplex: splitting the band or cancelling the echo}

A dial-up modem must send and receive at once over a two-wire loop. The early standards divided
the voice band: a V.21 modem originating a call transmits FSK around 1080\,Hz and receives around
1750\,Hz, and V.22bis used carriers of 1200 and 2400\,Hz, each direction occupying about half the
channel. This is simple and immune to echo, but it halves the usable bandwidth in each direction.

V.32 (1984) took the other route: both modems transmit across the whole band at the same time, and
each cancels the echo of its own transmission (Section~\ref{sec:ch12:echo}). The echo
canceller must handle two echoes: the \emph{near echo} from the local hybrid, strong and short,
and the \emph{far echo} from the distant four-wire-to-two-wire conversion, weaker but delayed by up
to a few hundred milliseconds on satellite routes, and subject to the frequency offset of analog
carrier systems along the way. With a cancellation of 50 to 60\,dB, both directions get the full
2400\,Bd channel. From V.32 onwards, every high-speed modem, ISDN, HDSL and Gigabit Ethernet has
cancelled echo rather than splitting the band.

\subsection{V.34: within a few decibels of capacity}

V.34 (1994, extended to 33.6\,kb/s in 1996) was the culmination of analog modem design. At start-up
the two modems \emph{probe the line}, sending a set of tones and measuring the received spectrum
and noise, and then choose a symbol rate (2400 to 3429\,Bd) and carrier frequency that best fit the
channel---a coarse, single-carrier version of the bit loading that DSL would do with many carriers.
Intersymbol interference is removed by Tomlinson--Harashima-style precoding at the transmitter
(Section~\ref{sec:ch12:thp}) rather than a DFE at the receiver, so that the trellis decoder sees
an ISI-free signal; four-dimensional trellis coding gives about 4 to 5\,dB of coding gain; and
\emph{shell mapping} shapes the constellation for close to 1\,dB of shaping gain
(Chapter~\ref{ch:modulation}). On a good line V.34 reaches about 93\% of the capacity computed in
Chapter~\ref{ch:infotheory}. Nothing more could be squeezed from an analog end-to-end voice channel.

\subsection{V.90: changing the channel model}

The 56\,kb/s modems of 1997--98 did not beat Shannon; they changed the channel. By the 1990s the
network between exchanges was entirely digital, so the only analog stage in a call from a home to
an Internet service provider was the subscriber's own loop. An ISP connected digitally to the
network (by T1 or ISDN primary-rate lines) could therefore send \emph{PCM codewords}, and the
exchange's codec at the subscriber's line card would convert each one to one of the 255 $\mu$-law
voltage levels on the loop. A V.90 client modem receives those levels at 8000 per second and, after
equalising the loop, decides which level was sent. The ``quantisation noise'' that limits V.34 has
vanished, because the transmitter chooses its symbols to be exactly the quantiser's levels
(Chapter~\ref{ch:infotheory} has the capacity argument).

The practical rate fell short of $8000\times7=56$\,kb/s for three reasons. The $\mu$-law levels near
zero are spaced only a fraction of a millivolt apart, too close to separate in the loop's noise,
so a V.90 modem uses a subset of perhaps 64 to 128 levels chosen during training. In North America,
robbed-bit signalling on some digital trunks (Chapter~\ref{ch:pcm}) corrupts the least significant
bit of every sixth sample. And regulators limited the average power on the line, which caps the use
of the large outer levels. Real connections ran at 40 to 53\,kb/s. The upstream direction still
crossed an analog-to-PCM conversion at the exchange and so stayed at V.34 rates until V.92 (2000)
applied precoding so that the exchange's codec would quantise the upstream signal onto intended
levels, reaching up to 48\,kb/s.

\begin{keyidea}[The channel is what you say it is]
V.34 treated the telephone network as a 3\,kHz analog channel with 35\,dB of SNR and reached its
capacity. V.90 treated the same network as a digital channel preceded by a known quantiser and a
short analog loop, and reached a capacity almost twice as high. The physics did not change; the
\emph{model} did. Many leaps in this chapter come from the same move: DSL used the loop's
frequencies above 4\,kHz that the voice network had thrown away, vectoring treated crosstalk as a
known MIMO channel rather than noise, and coherent optics treated the fibre as a linear field
channel rather than an intensity channel.
\end{keyidea}

%======================================================================
\section{The digital transmission hierarchy}
\label{sec:ch24:hierarchy}
%======================================================================

Chapter~\ref{ch:pcm} described how T1 and E1 carriers put 24 or 30 PCM voice channels on a pair of
wires and how these were multiplexed into higher rates. This section completes that story. It
explains why the first hierarchy was hard to manage, how SONET and SDH fixed it, and how the optical
transport network now carries everything, mostly Ethernet, over wavelengths.

\subsection{Plesiochronous multiplexing and bit stuffing}

In the plesiochronous digital hierarchy (PDH), each tributary is timed by its own clock, within a
tolerance of typically $\pm50$ parts per million for a DS1 and $\pm20$\,ppm for a DS3. A multiplexer
combining four DS1s into a DS2 cannot simply interleave their bits, because their rates differ
slightly and vary with time. Instead, it runs at a rate slightly higher than the sum of the inputs:
the DS2 at 6.312\,Mb/s carries $4\times1.544=6.176$\,Mb/s of tributaries plus 136\,kb/s of framing,
control and \emph{stuffing opportunities}. In each frame, each tributary has one bit position that
either carries a data bit or a dummy \emph{stuff} bit, according to whether that tributary's buffer
is running full or empty; \emph{stuffing control} bits, repeated several times and decided by
majority vote, tell the demultiplexer which. This is \term{positive justification}, and with it the
multiplexer can absorb any tributary rate within the tolerance. The demultiplexer removes the stuff
bits and smooths the gaps with a phase-locked loop, leaving a small residual \emph{waiting-time
jitter}, the slow wander in timing caused by the irregular pattern of stuffing events.

PDH worked, but it was hard to operate. Because every level stuffs independently, a DS1 is buried
at an unpredictable bit position inside a DS3 inside a 140\,Mb/s or 565\,Mb/s line; to drop or add a
single 1.5\,Mb/s circuit at an intermediate office, the operator had to demultiplex the whole
hierarchy level by level and multiplex it up again---the ``multiplexer mountain''. There was little
built-in performance monitoring, and the North American, European and Japanese hierarchies were
incompatible above the primary rate.

\subsection{SONET and SDH}

The \term{synchronous optical network} (SONET), developed at Bellcore and standardised by ANSI from
1988, and its international counterpart, the \term{synchronous digital hierarchy} (SDH, ITU-T G.707),
made the whole network synchronous to a hierarchy of atomic-referenced clocks and organised all rates
around a single frame repeating every 125\,$\mu$s.

\begin{figure}[htbp]\centering
\begin{tikzpicture}[x=0.11cm,y=0.42cm,font=\sffamily\scriptsize]
% TOH
\fill[accent!15] (0,0) rectangle (3,-3); \fill[accent!8] (0,-3) rectangle (3,-9);
\fill[navy!14] (3,0) rectangle (4,-9);
\fill[navy!5] (4,0) rectangle (90,-9);
\draw[navy,thick] (0,0) rectangle (90,-9);
\draw[navy] (3,0) -- (3,-9); \draw[navy] (4,0) -- (4,-9);
\draw[navy] (0,-3) -- (3,-3);
\foreach \r in {1,...,8} \draw[navy!25] (4,-\r) -- (90,-\r);
\node[rotate=90,align=center] at (1.5,-1.5) {SOH};
\node[rotate=90,align=center] at (1.5,-6) {LOH};
\node[rotate=90] at (3.5,-4.5) {POH};
\node[align=center] at (47,-4.5) {\small payload (86 columns $\times$ 9 rows = 774 bytes)\\[2pt]DS3, VT1.5s (DS1s), ATM cells, packets via GFP};
\draw[decorate,decoration={brace,amplitude=4pt},navy] (0,0.25) -- node[above=4pt]{3 col.} (3,0.25);
\draw[decorate,decoration={brace,amplitude=4pt},navy] (3,0.25) -- node[above=4pt]{synchronous payload envelope (87 columns), floats; located by the pointer} (90,0.25);
\draw[decorate,decoration={brace,amplitude=4pt,mirror},navy] (-0.6,0) -- node[left=4pt,align=right]{9 rows} (-0.6,-9);
\node[anchor=west,align=left,text=accent] at (91,-1.5) {framing A1A2,\\ parity B1, J0};
\node[anchor=west,align=left,text=accent] at (91,-4.2) {pointer H1H2H3,\\ parity B2,};
\node[anchor=west,align=left,text=accent] at (91,-6.6) {APS bytes K1K2,\\ DCC, sync S1};
\node[text=navy] at (45,-9.8) {transmitted row by row, left to right: 810 bytes every 125\,$\mu$s = 51.84\,Mb/s};
\end{tikzpicture}
\caption{The SONET STS-1 frame (SDH uses the same structure three times wider as the STM-1). The first
three columns are transport overhead: section overhead (framing, error monitoring, a data channel) and
line overhead (the payload pointer, protection-switching bytes, synchronisation status). The payload
envelope, with its own column of path overhead, can start anywhere in the frame; the pointer says where.}
\label{fig:ch24_sonet}
\end{figure}

The basic SONET signal, the STS-1 (synchronous transport signal level 1), is a frame of 9 rows of
90 bytes, sent row by row every 125\,$\mu$s (Figure~\ref{fig:ch24_sonet}): $9\times90\times8\times8000=
51.84$\,Mb/s. Higher rates are formed by byte-interleaving $N$ STS-1s into an STS-$N$, carried on
fibre as an OC-$N$ (optical carrier) at exactly $N\times51.84$\,Mb/s: OC-3 at 155.52\,Mb/s (equal to
SDH's STM-1), OC-12 at 622.08\,Mb/s, OC-48 at 2.48832\,Gb/s, OC-192 at 9.95328\,Gb/s and OC-768 at
39.81312\,Gb/s. A ``concatenated'' signal such as STS-3c treats the $N$ payloads as one large
container, for ATM cells or IP packets.

Three ideas made SONET a success.
\begin{description}[style=nextline,leftmargin=1.2em,font=\sffamily\bfseries\small\color{navy}]
\item[Pointers instead of stuffing.] The payload, the \emph{synchronous payload envelope} (SPE),
floats in the frame. A pointer in the H1 and H2 bytes of the line overhead gives the offset of the
start of the SPE. If the payload's clock is slightly fast, the multiplexer occasionally sends one
extra payload byte in the H3 byte (\emph{negative justification}) and decrements the pointer; if it
is slow, it leaves the byte after H3 empty (\emph{positive justification}) and increments the
pointer. The pointer moves in whole bytes, at most once every few frames, and any tributary, down to
a 1.5\,Mb/s \emph{virtual tributary} (VT1.5) carrying a DS1, can be located and extracted directly
without demultiplexing the rest. An \emph{add--drop multiplexer} can drop and insert individual
circuits at any node.
\item[Generous overhead.] Bit-interleaved parity bytes (B1, B2, B3) monitor errors on each section,
line and path; trace bytes verify connectivity; data communication channels carry management
traffic; and the K1 and K2 bytes carry the automatic protection switching protocol.
\item[Survivable rings.] SONET networks are built as rings of add--drop multiplexers, with every
circuit able to reach its destination either way round. In a \emph{unidirectional path-switched
ring} (UPSR, SDH's SNCP) every signal is sent both ways and the receiver selects the better copy; in
a \emph{bidirectional line-switched ring} (BLSR, SDH's MS-SPRing) half the capacity is reserved for
protection and, on a fibre cut, the two nodes adjacent to the break loop traffic back the other way.
Either way, service is restored in under 50\,ms (Telcordia GR-253, ITU-T G.841), short enough that
voice calls survive a cable cut with only a click. The ``50\,ms'' figure became a benchmark that
packet networks have been measured against ever since.
\end{description}

\begin{worked}[SONET overhead and payload efficiency]
How much of an STS-1 is available for payload, and how efficiently does it carry a DS3? How does
an OC-192c compare with 10\,Gigabit Ethernet?

Each 810-byte STS-1 frame spends $3\times9=27$\,bytes on transport overhead and 9 bytes on path
overhead, leaving $810-36=774$ bytes. At 8000 frames per second this is
$774\times64\,\text{kb/s}=49.536$\,Mb/s of payload capacity out of 51.84\,Mb/s: an overhead of only
4.4\%. A DS3 at 44.736\,Mb/s uses $44.736/51.84=86.3$\% of the line rate; the rest is SONET overhead
and fixed stuffing in the mapping.

An OC-192c (STS-192c) has 192 times the transport overhead but only one column of path overhead plus
63 columns of fixed stuff, leaving a payload capacity of 9.58464\,Gb/s, 96.3\% of the 9.95328\,Gb/s
line rate. Ethernet's designers took this number directly: the 10GBASE-W ``WAN PHY'' runs the 64b/66b
coded Ethernet stream at a payload rate of 9.58464\,Gb/s so that it fits an OC-192c, while the LAN
PHY 10GBASE-R runs at 10.3125\,GBd ($10\times66/64$). In practice operators preferred to carry
the LAN PHY transparently over OTN, the subject of the next subsection.
\end{worked}

\subsection{The optical transport network}

By 2000, most traffic was IP rather than voice, and each fibre carried many wavelengths. The
\term{optical transport network} (OTN, ITU-T G.709, first published in 2001) is the successor to
SONET/SDH for this world, a ``digital wrapper'' around any client signal on a wavelength. An OTN
frame has 4 rows of 4080 bytes: 16 columns of overhead, 3808 columns of payload (the optical payload
unit, OPU) and 256 columns of forward error correction. The standard FEC is the Reed--Solomon
RS(255,239) code, interleaved 16 ways, with 6.7\% overhead and about 6\,dB of net coding gain
(Chapter~\ref{ch:classiccodes}); stronger proprietary or standardised codes can replace it
(Chapter~\ref{ch:moderncodes}). The frame is the same at every rate, so its period shrinks as the
rate grows. The line rates are chosen to carry the corresponding SONET or Ethernet clients:
OTU1 at about 2.666\,Gb/s (for OC-48), OTU2 at 10.709\,Gb/s (OC-192), OTU3 at 43.018\,Gb/s
(OC-768) and OTU4 at 111.81\,Gb/s (100\,Gb Ethernet). Beyond 100\,Gb/s, the OTUC$n$ format
stacks $n$ instances of a roughly 105\,Gb/s frame, so that a 400 or 800\,Gb/s wavelength is simply
an OTUC4 or OTUC8, and the FlexO interfaces of ITU-T G.709.1 carry them between equipment.

OTN keeps SONET's virtues---overhead for monitoring each segment, protection switching, a
multiplexing hierarchy of optical data units (ODU0 at 1.25\,Gb/s up to ODUflex of any size)---and
adds strong FEC, which was essential as rates per wavelength rose and margins shrank. Most long-haul
traffic in the world travels in OTN frames, carrying Ethernet clients mapped through the generic
framing procedure (GFP) or directly.

%======================================================================
\section{Digital subscriber lines}
\label{sec:ch24:dsl}
%======================================================================

The subscriber loop's loss rises steeply with frequency (Section~\ref{sec:ch24:tline}), but at
1\,MHz a 3\,km loop still has a useful SNR. The voice network filtered off everything above
4\,kHz at the exchange; \term{digital subscriber line} (DSL) technology connects new equipment to the
bare copper at the exchange or at a street cabinet and uses the whole usable spectrum. The history
of DSL is a sequence of better and better answers to the question of how much data a given loop
can carry, and each answer is a chapter of this book in action.

\subsection{ISDN, HDSL and SHDSL}

The first digital service over the loop was the ISDN basic-rate interface (late 1980s), carrying two
64\,kb/s B channels and a 16\,kb/s D channel. In North America its line code (ANSI T1.601) is 2B1Q,
4-level PAM at 80\,kBd for a line rate of 160\,kb/s (Chapter~\ref{ch:baseband}); Germany and some
other countries used the ternary 4B3T code at 120\,kBd. Both directions share the full band, with
an echo canceller and an adaptive DFE at each end, over loops up to about 5.5\,km. ISDN set the
template of full-duplex baseband transmission over the loop.

\term{HDSL} (high-bit-rate DSL, early 1990s) was built to deliver T1 and E1 circuits to businesses
without repeaters. Instead of 1.544\,Mb/s AMI on two pairs with a repeater every 1.8\,km
(Chapter~\ref{ch:pcm}), HDSL sent 784\,kb/s of 2B1Q over each of two pairs, echo-cancelled, to the
full carrier-serving-area range of about 3.7\,km of 24-gauge cable. The single-pair successors,
HDSL2 and the international SHDSL (ITU-T G.991.2), use trellis-coded 16-level PAM (TC-PAM) with
Tomlinson--Harashima precoding and spectrally shaped transmit filters, carrying 192\,kb/s to
2.3\,Mb/s (and later about 5.7\,Mb/s with 32 levels) symmetrically. Symmetric business DSL of this
kind was the workhorse for leased lines and mobile backhaul in the 2000s.

\subsection{ADSL: asymmetric, multicarrier and over the voice}

\begin{history}[How DMT won the loop]
In 1989 Joseph Lechleider of Bellcore observed that the dominant impairment for a symmetric
high-speed service on the loop was self-NEXT, and that an \emph{asymmetric} service---fast towards
the subscriber, slow back---could put the two directions in different bands and escape it. The
motivating application was video on demand. At Stanford, John Cioffi's group showed that discrete
multitone modulation, an OFDM variant with bit loading (Chapter~\ref{ch:ofdm}), could approach the
capacity of the highly frequency-selective loop; Cioffi's company Amati built prototypes. In 1993
Bellcore ran a comparison of competing ADSL line codes, the ``ADSL Olympics'', in which Amati's DMT
outperformed the single-carrier carrierless amplitude and phase (CAP) modulation favoured by AT\&T
Paradyne. DMT was adopted in ANSI T1.413 (1995) and in ITU-T G.992.1 (1999). Video on demand did not
take off; the World Wide Web did, and its traffic was just as asymmetric. By the mid-2010s several
hundred million ADSL lines were in service, making it, for a time, the most widely deployed
broadband technology in the world.
\end{history}

ADSL (G.992.1) uses a DMT signal with 4.3125\,kHz tones: 256 downstream tones spanning 0 to
1.104\,MHz, generated by a 512-point real IFFT at 2.208\,MS/s. In the usual Annex~A variant over an
analog telephone line, tones 6 to 31 (25.875--138\,kHz) carry upstream data and tones 32 to 255
carry downstream data; a passive \emph{splitter} or in-home \emph{microfilters} separate the voice
band below 4\,kHz from the DSL band, so that the phone and the modem share the line. Annex~B, for
lines that also carry ISDN, moves upstream to 138--276\,kHz. The DMT symbol rate is 4000 per second
after a cyclic prefix and one synchronisation symbol per 68 data symbols; each tone carries up to 15
bits, determined during training from the measured SNR (Section~\ref{sec:ch17:coded}). A four-dimensional
16-state trellis code and an outer Reed--Solomon code with a convolutional interleaver protect the
data, the interleaver spreading the bursts of \emph{impulse noise} (from lightning, appliance motors and
ringing on neighbouring pairs) across many RS codewords. Its depth is a trade-off between impulse
protection and latency, and operators offered ``fast'' (non-interleaved) and ``interleaved'' paths for
gamers and for everyone else. Online reconfiguration---\emph{bit swapping}, which moves bits between
tones as the noise changes, and later \emph{seamless rate adaptation}---keeps the line up as conditions
drift.

ADSL2 (G.992.3, 2002) improved the framing, the initialisation and the diagnostics, and ADSL2+
(G.992.5, 2003) doubled the downstream band to 2.208\,MHz with 512 tones, raising the maximum rate
from about 8 to about 24\,Mb/s on short loops. Because loss rises so quickly with frequency, the
extra band helps only loops shorter than about 3\,km, and the rate of all forms of ADSL falls
steeply with distance (Figure~\ref{fig:ch24_dsl_rate_reach}).

\begin{worked}[The ADSL2+ rate on a 2\,km loop]
Estimate the downstream rate of ADSL2+ on a 2\,km, 0.4\,mm (26 AWG) loop with a transmit PSD of
$-40$\,dBm/Hz, background noise of $-140$\,dBm/Hz and FEXT from 10 other ADSL2+ lines in the binder.
Use an SNR gap of 12\,dB (9.8\,dB for uncoded QAM at a symbol error rate of $10^{-7}$, plus 6\,dB of
margin, less about 4\,dB of coding gain).

The loop model gives the loss at a few tones; the SNR is $-40-\text{loss}-N$, where $N$ is the noise
PSD: the $-140$\,dBm/Hz background plus the FEXT. Because FEXT travels down the same loop as the
signal, it dominates at the low frequencies, where it caps the SNR at about 60\,dB, and fades into
the background noise towards the top of the band:
\begin{center}\small
\begin{tabular}{@{}lcccc@{}}
\toprule
Frequency & 0.2\,MHz & 0.5\,MHz & 1.0\,MHz & 1.5\,MHz\\
\midrule
Loop loss (dB) & 25 & 36 & 51 & 63\\
SNR (dB) & 61 & 53 & 45 & 36\\
$b=\lfloor\log_2(1+\text{SNR}/\Gamma)\rfloor$ & 15 (cap) & 13 & 10 & 8\\
\bottomrule
\end{tabular}
\end{center}
At 1\,MHz, for example, SNR$/\Gamma=45-12=33$\,dB, a factor of 2000, and $\log_2 2001=10.97$, so 10
bits. The loading falls by roughly one bit per 100\,kHz above 0.5\,MHz, and every tone up to 2.2\,MHz
still carries a few bits. Summing over tones 33 to 511 gives about 4470 bits per DMT symbol,
or $4470\times4000\approx17.9$\,Mb/s of line rate. At 3\,km the same calculation gives about 8\,Mb/s
with no usable tones above 1.3\,MHz, and at 4\,km about 3.6\,Mb/s. Chapter~\ref{ch:ofdm} does the
same calculation with a harsher noise model including an AM radio interferer and obtains lower
rates: real rates depend strongly on the noise environment, which is why operators quote
rate-versus-reach curves for several noise scenarios. These are line rates; the net rate after
framing and FEC overhead is roughly 10--15\% lower.
\end{worked}

\mdcfig[0.85]{ch24_dsl_rate_reach}{Rate versus reach of the main DSL families on 0.4\,mm (26 AWG)
cable, computed by gap-approximation bit loading (12\,dB gap, $-140$\,dBm/Hz background noise, Unger
FEXT models with 10 ADSL or 20 VDSL2/G.fast disturbers). ADSL is capped at about 8\,Mb/s by its framing.
VDSL2 profile 17a (downstream bands of the 998ADE17 plan, 14.5\,dBm aggregate power) gains little
over ADSL2+ beyond about 1.5\,km, but vectoring adds substantially on short loops. G.fast (106\,MHz,
time-division duplex, shown as aggregate down plus up rate) reaches about 1\,Gb/s only on loops of
tens of metres. Each step up in rate requires the fibre to come closer to the home.}

\subsection{VDSL2 and the fibre-fed cabinet}

Figure~\ref{fig:ch24_dsl_rate_reach} makes the economics plain: the only way to get more rate from
copper is to make the copper shorter. In the 2000s operators began to install DSL equipment in
street cabinets fed by fibre (\emph{fibre to the cabinet}, FTTC, or \emph{to the node}, FTTN), leaving
loops of a few hundred metres to a kilometre. \term{VDSL2} (ITU-T G.993.2, 2006) was designed for
them. It keeps ADSL's DMT, framing and coding but extends the band: profile 8 up to 8.5\,MHz,
profile 12 to 12\,MHz, profile 17a to 17.664\,MHz with 4096 tones, and profile 30a to 30\,MHz with
8.625\,kHz tones; the later profile 35b (2015) uses 35\,MHz. The bands alternate between downstream
and upstream according to a band plan (in Europe, the ``998'' plan with downstream bands up to 3.75,
5.2--8.5 and 12--17.664\,MHz), so that each direction has some low-loss spectrum on short loops and
the plan works across loop lengths. Transmit power is limited in aggregate (14.5\,dBm downstream for
profile 17a) as well as by a PSD mask, so the transmitter spreads its power over the usable tones.

\begin{figure}[htbp]\centering
\begin{tikzpicture}[font=\sffamily\scriptsize,
  box/.style={draw=navy,fill=navy!6,thick,rounded corners=2pt,minimum height=0.6cm,minimum width=1.25cm,align=center},
  fib/.style={very thick,accent},
  cop/.style={very thick,navy!70,decorate,decoration={snake,amplitude=0.7pt,segment length=4pt}},
  home/.style={draw=navy,fill=white,thick,minimum height=0.45cm,minimum width=0.75cm,align=center}]
\node[box,minimum height=6.0cm,minimum width=1.5cm] (co) at (0,-0.45) {Central\\office\\[2pt](OLT,\\DSLAM,\\BNG)};
\foreach \y/\lab in {2.0/FTTEx, 0.6/FTTC, -0.8/FTTdp, -2.3/FTTH} {\node[anchor=east,text=navy] at (-0.9,\y) {\textbf{\lab}};}
% FTTEx
\draw[cop] (0.75,2.0) -- (9.3,2.0); \node[home] at (9.75,2.0) {home};
\node[above,text=navy] at (5.0,2.1) {copper loop up to about 5\,km: ADSL2+ (G.992.5), 1--24\,Mb/s};
% FTTC
\draw[fib] (0.75,0.6) -- (6.4,0.6); \node[box] (cab) at (7.05,0.6) {cabinet\\VDSL2};
\draw[cop] (7.7,0.6) -- (9.3,0.6); \node[home] at (9.75,0.6) {home};
\node[above,text=navy] at (3.6,0.65) {fibre to the cabinet};
\node[above,text=navy] at (8.5,0.7) {0.3--1\,km};
% FTTdp
\draw[fib] (0.75,-0.8) -- (7.7,-0.8); \node[box] (dp) at (8.35,-0.8) {DPU\\G.fast};
\draw[cop] (9.0,-0.8) -- (9.3,-0.8); \node[home] at (9.75,-0.8) {home};
\node[above,text=navy] at (4.2,-0.75) {fibre to the distribution point};
\node[below,text=navy] at (8.35,-1.15) {copper $<$250\,m};
% FTTH
\draw[fib] (0.75,-2.3) -- (6.2,-2.3);
\node[draw=accent,fill=accent!6,thick,isosceles triangle,isosceles triangle apex angle=60,inner sep=1pt,minimum size=0.55cm] (spl) at (6.5,-2.3) {};
\node[above=0.12cm of spl,text=accent] {1:32 splitter};
\foreach \dy in {0.4,0,-0.4} \draw[fib] (spl.east) -- (9.3,-2.3+\dy);
\node[home] at (9.75,-1.9) {ONT}; \node[home] at (9.75,-2.3) {ONT}; \node[home] at (9.75,-2.7) {ONT};
\node[above,text=navy] at (3.4,-2.25) {feeder fibre, up to 20\,km};
\node[below,text=navy] at (3.4,-2.4) {PON: GPON, XGS-PON, 50G-PON};
\draw[fib] (2.0,-3.55) -- (2.8,-3.55) node[right,text=navy]{fibre};
\draw[cop] (4.4,-3.55) -- (5.2,-3.55) node[right,text=navy]{copper pair};
\end{tikzpicture}
\caption{Fibre-to-the-x access architectures. Each step pushes the fibre closer to the subscriber and
shortens the copper, raising the rate: ADSL2+ from the exchange, VDSL2 from a street cabinet, G.fast
from a distribution point within a few hundred metres, and finally a passive optical network all the
way to an optical network terminal (ONT) in the home.}
\label{fig:ch24_fttx}
\end{figure}

\subsection{Vectoring: crosstalk cancellation as MIMO}
\label{sec:ch24:vectoring}

On the short loops of VDSL2, the received SNR on each tone is set not by background noise but by
FEXT from the other VDSL2 lines in the binder (Figure~\ref{fig:ch24_crosstalk_vectoring}). FEXT is not
noise, however: it is a deterministic, slowly varying linear coupling of signals that the DSLAM
itself transmits. Consider tone $k$ on the $N$ lines served by one DSLAM. The received samples are
\begin{equation}
\vect y_k=\vect H_k\vect x_k+\vect n_k,
\label{eq:ch24:vectmodel}
\end{equation}
where $\vect x_k$ and $\vect y_k$ are the $N$ transmitted and received values on that tone and
$\vect H_k$ is an $N\times N$ channel matrix whose diagonal holds the direct (through) channels and
whose off-diagonal entries are the FEXT couplings. This is exactly the MIMO model of
Chapter~\ref{ch:mimo}, one copy per tone, with one important simplification: the matrix is strongly
\emph{diagonally dominant}, because each FEXT coupling is 20 to 60\,dB below the direct channel.
Write $\vect H_k=\vect D_k(\vect I+\vect C_k)$, with $\vect D_k$ diagonal and the crosstalk
matrix $\vect C_k$ small.

\begin{description}[style=nextline,leftmargin=1.2em,font=\sffamily\bfseries\small\color{navy}]
\item[Upstream: joint cancellation.] The receivers of all $N$ lines are in the same DSLAM, so it can
process $\vect y_k$ jointly. A zero-forcing canceller computes $(\vect I+\vect C_k)^{-1}\vect D_k^{-1}\vect y_k$
or, to first order, subtracts $\vect C_k$ times the decisions of the other lines, a multiuser
generalisation of the DFE. Because $\vect H_k$ is close to diagonal, zero forcing loses almost nothing
to noise enhancement.
\item[Downstream: precoding.] The receivers are in different homes and cannot cooperate, but the
transmitters are together. The DSLAM therefore \emph{precodes}, transmitting
$\vect x_k=\vect P_k\vect u_k$ with $\vect P_k=(\vect I+\vect C_k)^{-1}\approx\vect I-\vect C_k$, so that the
received signal is $\vect y_k=\vect D_k\vect u_k+\vect n_k$: each line receives its own data with the
crosstalk pre-cancelled by deliberately added anti-crosstalk. Since $\vect C_k$ is small, the
precoder barely raises the transmit power, which the PSD mask requires.
\end{description}

The crosstalk channel must be estimated, and that is the clever part of the vectoring standard,
ITU-T G.993.5 (2010), which builds on the analysis of Ginis and Cioffi (2002). During the periodic
synchronisation symbols, every line transmits a known pilot sequence, the sequences of different
lines being mutually orthogonal (Walsh--Hadamard sequences, as in CDMA, Chapter~\ref{ch:spreadspectrum}).
Each customer modem reports back its received error samples on those symbols; correlating the errors
against each line's pilot sequence gives the corresponding column of $\vect C_k$. Precoding can thus
be added to the downstream of existing customer modems with only a firmware change in the modem to
report errors.

\mdcfig{ch24_crosstalk_vectoring}{Left: on 1\,km of 26 AWG cable, NEXT from 10 disturbers exceeds the
through signal above about 2\,MHz, which is why DSL separates the directions in frequency or time;
FEXT is weaker but rises as $f^2$ relative to the signal. Right: bit loading of VDSL2 profile 17a
downstream on a 600\,m loop with 20 same-service disturbers. Without vectoring, FEXT costs about a
quarter of the rate; cancelling the crosstalk by 25\,dB recovers essentially all of it, approaching
the crosstalk-free bound (wide grey trace).}

With vectoring, the rate of a VDSL2 line approaches its single-line, crosstalk-free value: our model
gives about 69\,Mb/s without and 91\,Mb/s with vectoring on a 600\,m loop, and the gain is larger on
shorter loops (about 92 versus 160\,Mb/s at 300\,m), where FEXT matters most. The computational load
is large: a DSLAM vectoring 192 lines on 4096 tones at 4000 symbols per second performs on the order
of $192^2\times4096\times4000\approx6\times10^{11}$ complex multiply--accumulates per second in each
direction, which is why practical systems limit the number of crosstalkers cancelled per line and per
tone (\emph{partial vectoring}).

\begin{pitfall}[Vectoring needs every line in the binder]
Crosstalk from a line that is \emph{not} part of the vectored group, called \emph{alien crosstalk},
cannot be cancelled, because its signal is unknown to the precoder. A single non-vectored VDSL2 line in a
binder can erase much of the vectoring gain for all its neighbours. That turned a technical constraint
into a regulatory one: where several operators had used unbundled loops in the same cable, regulators
had to decide whether one operator could control all VDSL2 lines in a cabinet, or require
multi-operator ``node-level'' vectoring schemes. When you plan a vectored deployment, inventory every
signal in the binder.
\end{pitfall}

\subsection{G.fast and the distribution point}
\label{sec:ch24:gfast}

\term{G.fast} (ITU-T G.9701, 2014) pushes the same logic to its limit: fibre to a \emph{distribution
point unit} (DPU) on a pole, in a manhole or in a building's basement, serving perhaps 8 to 48
subscribers over loops shorter than about 250\,m. To use frequencies up to 106\,MHz (later
212\,MHz), it makes several changes:
\begin{itemize}
\item \emph{Time-division duplexing} with synchronised frames: all lines in a binder transmit
downstream at the same time and upstream at the same time, so there is no NEXT, and the
downstream/upstream split of the frame can be set by the operator. The DMT tone spacing is
51.75\,kHz (twelve times ADSL's) with 2048 tones for 106\,MHz, and the symbol rate is 48\,000 per second.
\item \emph{Vectoring is built in from the start}, and at these frequencies the FEXT can be as strong as
the direct signal, so the first-order precoder $\vect I-\vect C$ is not enough. G.fast allows full
matrix inversion and \emph{nonlinear precoding} with a modulo operation, the multiuser form of
Tomlinson--Harashima precoding (Section~\ref{sec:ch12:thp}), which avoids the power increase of linear
precoding when the channel is far from diagonal.
\item \emph{Reverse power feeding} (ETSI TS 101 548): the DPU, which may have no mains supply, is
powered over the copper pairs from the customers' homes.
\end{itemize}
On loops of 50 to 100\,m G.fast delivers on the order of 1\,Gb/s aggregate in its 106\,MHz profile,
and about twice that in the 212\,MHz profile; our simple model in Figure~\ref{fig:ch24_dsl_rate_reach}
gives 1.08\,Gb/s at 50\,m and 0.45\,Gb/s at 200\,m. Its successor, MGfast (ITU-T G.9711), extends the
band further, to several hundred megahertz, for in-building runs. These are probably the last
generations of copper access; the next step, shown in Figure~\ref{fig:ch24_fttx}, is fibre all the way.

\begin{inpractice}[DSL by the numbers]
Typical service figures, which depend strongly on the plant: ADSL2+ from an exchange, 8 to 20\,Mb/s
within 2\,km, a few megabits per second at 4\,km; VDSL2 17a vectored from a cabinet, about 100\,Mb/s
at 400 to 500\,m on 0.5\,mm cable; VDSL2 35b, up to about 300\,Mb/s on loops of a few hundred metres;
G.fast 106a, 500\,Mb/s to 1\,Gb/s aggregate within 100\,m. Cable gauge matters a great deal: 0.5\,mm
cable reaches roughly 25 to 30\% farther than 0.4\,mm for the same rate. Plant quality matters even
more: bridged taps, corroded joints and unbalanced pairs (which convert common-mode RF interference
into differential noise) are the usual causes of a line that underperforms its length, and DSL
modems report per-tone SNR and \emph{Hlog} (the measured channel magnitude) precisely so that such
faults can be diagnosed remotely.
\end{inpractice}

%======================================================================
\section{Cable networks: HFC and DOCSIS}
\label{sec:ch24:cable}
%======================================================================

The second copper network that reaches most homes in North America, and many elsewhere, was built
for television. Community antenna television (CATV) began in the late 1940s in mountain towns of
Pennsylvania and Oregon, where a shared antenna on a hilltop fed a coaxial cable to houses in the
valley. By the 1980s cable systems were trees of 75\,$\Omega$ coax with cascades of 20 to 40 broadband
amplifiers, carrying dozens of 6\,MHz analog TV channels (8\,MHz in Europe) one way, from the
\emph{headend} to the subscribers.

\subsection{Hybrid fibre--coax}

Long amplifier cascades accumulate noise and distortion and fail often. In the 1990s operators
replaced the trunk part of the tree with fibre carrying the whole RF spectrum as an analog intensity
modulation of a laser, terminating at an \emph{optical node} that converts it back to RF for a short
coaxial distribution network (Figure~\ref{fig:ch24_hfc}). This \term{hybrid fibre--coax} (HFC)
architecture shortened cascades to a few amplifiers, and adding a return path in the low part of the
spectrum (5--42\,MHz in North America, 5--65\,MHz in Europe) made the network two-way. The number of
homes per node has fallen steadily, from about 2000 in the 1990s to a few hundred, and the newest
designs are ``node plus zero'' (N+0), with no amplifiers at all after the node.

\begin{figure}[htbp]\centering
\begin{tikzpicture}[font=\sffamily\scriptsize,
  box/.style={draw=navy,fill=navy!6,thick,rounded corners=2pt,minimum height=0.7cm,align=center},
  amp/.style={draw=navy,fill=navy!10,thick,isosceles triangle,isosceles triangle apex angle=60,inner sep=0.5pt,minimum size=0.42cm},
  tap/.style={draw=navy,fill=white,thick,minimum size=0.22cm,inner sep=0pt},
  fib/.style={very thick,accent},
  cx/.style={thick,navy},
  home/.style={draw=navy,fill=white,thick,minimum height=0.36cm,minimum width=0.55cm,font=\sffamily\tiny}]
\node[box,minimum height=2.0cm,minimum width=1.9cm] (he) at (0,0) {Headend / hub\\[2pt]CMTS\\video\\(or Remote PHY\\core)};
\node[box,fill=accent!6,draw=accent] (node) at (4.2,0) {optical\\node};
\draw[fib] (he.east) -- node[above]{fibre, 10--40\,km} node[below,align=center]{analog RF on light, or\\digital (Remote PHY)} (node.west);
% main coax
\draw[cx] (node.east) -- (12.2,0);
\node[amp] at (7.2,0) {}; \node[amp] at (10.4,0) {};
\foreach \x in {5.6,6.4,8.2,9.0,9.8,11.2,12.0} {\node[tap] (t\x) at (\x,0) {}; \draw[cx,thin] (\x,0) -- (\x,-0.85); \node[home] at (\x,-1.1) {CM};}
\draw[cx] (8.6,0) -- (8.6,1.0) -- (11.5,1.0); \node[amp] at (9.4,1.0) {};
\foreach \x in {10.2,11.0} {\node[tap] at (\x,1.0) {}; \draw[cx,thin] (\x,1.0) -- (\x,1.6); \node[home] at (\x,1.85) {CM};}
\node[text=navy] at (8.6,-1.65) {taps and drops to cable modems (CM); 75\,$\Omega$ coax};
\node[text=navy] at (7.2,0.42) {amplifier};
\draw[fib] (0.1,-2.2) -- (0.9,-2.2) node[right,text=navy]{fibre};
\draw[cx] (2.2,-2.2) -- (3.0,-2.2) node[right,text=navy]{coax};
\end{tikzpicture}
\caption{A hybrid fibre--coax network. The headend's cable modem termination system (CMTS) and video
equipment drive optical links to nodes, each serving a few hundred homes over a tree of coaxial cable
with amplifiers and taps. In a distributed access architecture the PHY moves into the node (Remote PHY)
and the fibre carries Ethernet instead of analog RF.}
\label{fig:ch24_hfc}
\end{figure}

The coaxial plant is a channel quite different from a twisted pair. Its loss is lower but still rises
as $\sqrt f$, so amplifiers are set with an upward \emph{tilt}, launching higher frequencies at
higher levels. Its echoes come from impedance mismatches at taps, splitters and connectors, typically
a few hundred nanoseconds to a few microseconds long. And its upstream is shared: the noise and ingress
picked up by every home on a node---shortwave radio, appliance impulses, and the infamous loose
connector---are added together at the node, a phenomenon cable engineers call \emph{noise funnelling},
which makes the upstream far noisier than the downstream.

\subsection{DOCSIS from 1.0 to 3.0}

The CableLabs \term{DOCSIS} specifications (Data Over Cable Service Interface Specification,
standardised internationally as ITU-T J.112, J.122 and successors) define the cable modem (CM) and the
\emph{cable modem termination system} (CMTS) at the headend. DOCSIS~1.0 (1997) used one 6\,MHz
television channel for downstream data with 64- or 256-QAM and the concatenated Reed--Solomon and
trellis coding of ITU-T J.83 Annex~B: 256-QAM at 5.36\,MBd gives about 38\,Mb/s of payload. The
upstream used bursts of QPSK or 16-QAM in channels up to 3.2\,MHz wide, shared by time-division
multiple access: modems request transmission opportunities, and the CMTS broadcasts \emph{MAP}
messages granting mini-slots. Each modem is \emph{ranged} at start-up---its timing advance, power and
frequency adjusted by the CMTS---so that bursts from all modems arrive aligned, exactly the
timing-advance problem of cellular uplinks (Chapter~\ref{ch:multipleaccess}). DOCSIS~2.0 (2001) widened
the upstream to 6.4\,MHz with 64-QAM, and DOCSIS~3.0 (2006) introduced \emph{channel bonding}, spreading
one modem's traffic over many channels: 32 downstream and 8 upstream channels were typical of later 3.0
modems, for roughly 1\,Gb/s downstream.

\subsection{DOCSIS 3.1: OFDM, LDPC and 4096-QAM}

DOCSIS~3.1 (2013) replaced the single-carrier channels with OFDM (Chapter~\ref{ch:ofdm}). A downstream
OFDM channel is up to 192\,MHz wide, with 4096- or 8192-point FFTs at 204.8\,MS/s, giving subcarrier
spacings of 50 or 25\,kHz, and a cyclic prefix selectable from 0.94 to 5\,$\mu$s to cover the plant's
echoes. The FEC is an LDPC code, (16\,200,14\,400) downstream with an outer BCH code, decoded with soft
decisions (Chapter~\ref{ch:moderncodes}), which gains several decibels over the J.83 code and allows
4096-QAM downstream, with 8192- and 16384-QAM optional. The upstream uses OFDMA in channels up to
96\,MHz wide. Rather than bit-loading every subcarrier as DSL does, DOCSIS~3.1 defines a small number
of \emph{modulation profiles}, each assigning constellations to ranges of subcarriers, and assigns each
modem the most efficient profile its measured modulation error ratio supports. The spectrum grew too:
upstream up to 204\,MHz and downstream up to 1218\,MHz (optionally 1794\,MHz), for about 10\,Gb/s
downstream and 1 to 2\,Gb/s upstream per node.

\mdcfig{ch24_docsis_spectrum}{Spectrum plans of DOCSIS generations (North American values; other
splits are allowed). DOCSIS 3.0 bonds 6\,MHz single-carrier QAM channels; DOCSIS 3.1 adds OFDM blocks
up to 192\,MHz wide and a high upstream split; DOCSIS 4.0 either shares 108--684\,MHz between the two
directions with echo cancellation (full duplex) or extends the spectrum to 1.8\,GHz with an upstream
split as high as 684\,MHz (extended-spectrum FDD).}

\begin{worked}[The capacity of a DOCSIS 3.1 OFDM channel]
A 192\,MHz downstream channel uses the 4K FFT (50\,kHz subcarriers, useful symbol time
$1/50\,\text{kHz}=20\,\mu$s) with a 1.25\,$\mu$s cyclic prefix. Estimate its throughput with 4096-QAM.

The 192\,MHz channel holds 3840 subcarriers, of which about 3800 remain after the guard bands at
the edges, and roughly 3700 after pilots and the PHY link channel. The OFDM symbol lasts
$20+1.25=21.25\,\mu$s. With 12 bits per subcarrier the raw rate is
$3700\times12/21.25\,\mu\text{s}\approx2.09$\,Gb/s. The LDPC code has rate $14\,400/16\,200=0.89$, and
the outer BCH code removes a little more, leaving about 1.85\,Gb/s of payload per 192\,MHz, or about
9.6\,b/s/Hz. Five such channels between 258 and 1218\,MHz give roughly the 10\,Gb/s headline figure.
Note how much this depends on the plant: each step down in constellation (to 2048- or 1024-QAM) costs
about 8\% of the rate, and requires 3\,dB less SNR---which is why operators measure and maintain
the plant's MER so carefully.
\end{worked}

\subsection{DOCSIS 4.0: full duplex and extended spectrum}

DOCSIS~4.0 (CableLabs, about 2020) targets 10\,Gb/s downstream and 6\,Gb/s upstream, and offers two
options. \emph{Extended-spectrum DOCSIS} keeps frequency-division duplexing but pushes the top of the
spectrum to 1794\,MHz and lets operators move the upstream/downstream split as high as 684\,MHz,
which requires new amplifiers and taps. \emph{Full-duplex DOCSIS} (FDX) keeps the plant at 1.2\,GHz but
uses the band from 108 to 684\,MHz in \emph{both} directions at once. The node, now with no amplifiers
after it, cancels the echo of its own downstream transmission from the upstream it receives, the same
echo cancellation as V.32 and Gigabit Ethernet at hundreds of megahertz of bandwidth. Modems cannot
cancel the upstream transmissions of their neighbours, which leak into their downstream through the
taps; the CMTS therefore measures which modems interfere with one another and groups them into
\emph{interference groups} that are never scheduled to transmit upstream in a sub-band while another
member receives downstream in it. In parallel, the industry moved to \emph{distributed access
architectures}: the PHY (Remote PHY) or the whole MAC moves into the node, and the fibre from the hub
carries digital Ethernet rather than analog RF, eliminating the noise and distortion of analog optics.

\begin{inpractice}[Cable versus fibre]
A DOCSIS~4.0 node delivers 10\,Gb/s shared by a few hundred homes, comparable to an XGS-PON
(Section~\ref{sec:ch24:pon}) shared by 32 or 64. Cable's advantage is that the coax is already there
and the upgrade is in the electronics; fibre's is that its capacity is limited by the optics, not the
medium, and that a passive network has no amplifiers to power and maintain. Many cable operators run
both: DOCSIS on the existing plant and PON for new builds, sometimes from the same hub under the same
management system.
\end{inpractice}

%======================================================================
\section{Ethernet over copper}
\label{sec:ch24:ethernet}
%======================================================================

Ethernet began on coaxial cable (Chapter~\ref{ch:baseband} tells why it used Manchester coding) and
moved in 1990 to the cheap twisted-pair wiring already installed in office buildings. Its
twisted-pair physical layers are a parade of the techniques of this book, each generation reaching
roughly ten times the rate over the same 100\,m of cable (Table~\ref{tab:ch24:ethernet}).

\begin{table}[htbp]\centering\small
\caption{Twisted-pair Ethernet physical layers (IEEE 802.3; years are approximate approval dates).}
\label{tab:ch24:ethernet}
\begin{tabularx}{\linewidth}{@{}l l l l X@{}}
\toprule
PHY & Clause/project & Year & Medium & Signalling\\
\midrule
10BASE-T & 802.3i & 1990 & 2 pairs Cat 3, 100\,m & Manchester, 20\,MBd\\
100BASE-TX & 802.3u & 1995 & 2 pairs Cat 5, 100\,m & 4B5B, scrambler, MLT-3 at 125\,MBd\\
1000BASE-T & 802.3ab & 1999 & 4 pairs Cat 5, 100\,m & PAM-5 at 125\,MBd per pair, 4D 8-state TCM, echo and NEXT cancellation\\
10GBASE-T & 802.3an & 2006 & 4 pairs Cat 6A, 100\,m & PAM-16 (DSQ128) at 800\,MBd, LDPC, THP\\
2.5G/5GBASE-T & 802.3bz & 2016 & 4 pairs Cat 5e/6, 100\,m & 10GBASE-T scaled to 1/4 and 1/2 the symbol rate\\
25G/40GBASE-T & 802.3bq & 2016 & 4 pairs Cat 8, 30\,m & 10GBASE-T scaled up\\
100BASE-T1 & 802.3bw & 2015 & 1 pair, 15\,m (automotive) & PAM-3 at 66.67\,MBd, full duplex\\
1000BASE-T1 & 802.3bp & 2016 & 1 pair, 15--40\,m & PAM-3 at 750\,MBd, Reed--Solomon FEC\\
10BASE-T1S & 802.3cg & 2019 & 1 pair, multidrop & Short-reach shared bus for sensors\\
\bottomrule
\end{tabularx}
\end{table}

\subsection{10 and 100\,Mb/s}

10BASE-T (1990) sends 10\,Mb/s Manchester code at 20\,MBd over one pair in each direction, with link
pulses to signal that a link is alive. 100BASE-TX (1995) could not afford Manchester's doubled
bandwidth on Category~5 cable, and instead used 4B5B coding at 125\,MBd, scrambled to spread the
spectrum and shaped by the three-level MLT-3 code so that most of the energy lies below about 30\,MHz
(Figure~\ref{fig:ch24_ethernet}). Both use one pair for each direction, so there is no echo, and the
receiver needs only a simple adaptive equaliser for the cable's $\sqrt f$ loss.

\subsection{1000BASE-T: four pairs, both directions, five levels}

Gigabit Ethernet over copper (IEEE 802.3ab, 1999) had to deliver ten times the rate over the same
Category~5 cable, specified only to 100\,MHz. The design, a landmark of DSP-based transceivers,
combined four ideas.
\begin{enumerate}
\item \emph{Use all four pairs}, each carrying 250\,Mb/s, at the same 125\,MBd symbol rate as
100BASE-TX, so the spectrum stays within the cable's specification.
\item \emph{Send two bits per symbol} with five-level PAM ($\{-2,-1,0,1,2\}$), the extra level
providing redundancy for a code.
\item \emph{Code across the pairs}: the four PAM-5 symbols sent at the same instant form a
four-dimensional point, protected by an 8-state four-dimensional trellis code in the spirit of
Ungerboeck's TCM (Chapter~\ref{ch:modulation}), whose coding gain more than repays the SNR cost of the
fifth level. The receiver combines Viterbi decoding with decision feedback in a decision-feedback
sequence estimator.
\item \emph{Transmit in both directions on every pair simultaneously}, with a hybrid at each end.
\end{enumerate}
The last idea is what made the design hard. Each receiver hears the far-end signal attenuated by the
cable, its own transmitter's echo from the hybrid and the cable's impedance discontinuities, and the
near-end crosstalk from the three other local transmitters, as well as far-end crosstalk
(Figure~\ref{fig:ch24_ethernet}, right). At the upper end of the band the echo is stronger than the signal
and the NEXT is not far below it. Each of the four receivers therefore runs an adaptive echo canceller and three NEXT
cancellers---FIR filters of tens to over a hundred taps fed by the local transmit symbols
(Section~\ref{sec:ch12:echo})---plus an FFE and DFE for the channel. Because the cancellers need the
transmit and receive clocks to be locked, one end is the timing \emph{master} and the other recovers
the master's clock and uses it for its own transmitter (\emph{loop timing}).

\mdcfig{ch24_ethernet}{Left: spectra of successive twisted-pair Ethernet PHYs (normalised): Manchester
10BASE-T, MLT-3 100BASE-TX, PAM-5 1000BASE-T (with its $0.75+0.25D$ transmit shaping) and PAM-16
10GBASE-T at 800\,MBd. Right: approximate levels at a 1000BASE-T receiver on a worst-case 100\,m
Category~5e channel, from the TIA-568 limit formulas: the far-end signal falls with the cable's loss,
while the echo (limited by the channel return loss) and the power-sum NEXT from the three other pairs
rise to approach it; above about 40\,MHz the echo is the strongest signal at the receiver. Without
cancellation the link could not work.}

\begin{worked}[How much echo cancellation does 1000BASE-T need?]
At 62.5\,MHz, the Nyquist frequency of 125\,MBd, a worst-case 100\,m Category~5e channel loses about
18\,dB, and its return loss may be as poor as about 12\,dB. Suppose the decoder needs an SNR of
about 22\,dB at the slicer (a rough figure for PAM-5 with the trellis code and some margin).

Relative to the transmitted level, the signal is at $-18$\,dB and the echo at $-12$\,dB, so the echo
is 6\,dB \emph{stronger} than the signal. For the residual echo to sit, say, 28\,dB below the signal
(so that, together with NEXT, FEXT and noise, the total stays 22\,dB below), the canceller must
suppress the echo by $6+28=34$\,dB, which requires an accurate model of an echo response tens of
symbols long, adapted continuously. NEXT, at about $-30$\,dB at 62.5\,MHz, only 12\,dB below the signal, must be
cancelled by roughly 16 to 20\,dB. In 1999 this was a demanding piece of mixed-signal DSP; within a few years it
cost cents per port.
\end{worked}

\subsection{10GBASE-T and its descendants}

10GBASE-T (802.3an, 2006) raised the symbol rate to 800\,MBd per pair over Category~6A cable
specified to 500\,MHz. Each pair carries 16-level PAM; pairs of consecutive symbols form a
two-dimensional point of the 128-point \emph{DSQ128} constellation (a checkerboard subset of the
$16\times16$ grid, 3.5 bits per symbol), and a (2048,1723) LDPC code protects the bits that select
between closely spaced points, while the others ride uncoded in the manner of multilevel coding.
Decision feedback equalisation is incompatible with an LDPC decoder whose decisions arrive thousands
of symbols late, so the channel's ISI is pre-equalised at the transmitter by Tomlinson--Harashima
precoding (Section~\ref{sec:ch12:thp}), with coefficients computed by the receiver during training and
sent back. Echo and NEXT cancellers with hundreds of taps each run at 800\,MBd. Early 10GBASE-T PHYs
dissipated several watts per port; process scaling brought that to a few watts, and the
\emph{multi-gigabit} 2.5GBASE-T and 5GBASE-T (802.3bz, 2016), the same design at a quarter and half the
symbol rate, now connect Wi-Fi~6 and~7 access points over existing Category~5e and~6 cable.

\subsection{Single-pair, backplane and powered Ethernet}

Cars, factories and buildings need Ethernet on a single pair, light, cheap and robust against harsh
electromagnetic environments. \emph{Automotive Ethernet} began with 100BASE-T1 (802.3bw, 2015), derived
from Broadcom's BroadR-Reach: full duplex on one unshielded pair with echo cancellation and PAM-3 at
66.67\,MBd, using the 3B2T mapping of Chapter~\ref{ch:baseband}. 1000BASE-T1 (802.3bp, 2016) runs
PAM-3 at 750\,MBd with Reed--Solomon FEC, and the multi-gigabit 802.3ch (2020) reaches 10\,Gb/s for
camera and sensor links; 10BASE-T1S (802.3cg) is a multidrop bus that replaces CAN-type networks.

Inside equipment, the \emph{backplane} and twinaxial \emph{direct-attach copper} (DAC) PHYs are the
SerDes links of Chapter~\ref{ch:baseband}: 10GBASE-KR (802.3ap, 2007) with NRZ at 10.3125\,GBd, and
100\,Gb/s-per-lane PAM-4 at 53.125\,GBd in 802.3ck, with 200\,Gb/s per lane following in 802.3dj. They
use the same FFE/CTLE/DFE equalisers and Reed--Solomon FEC described there.

Finally, \term{Power over Ethernet} (PoE) delivers DC power over the data pairs, using the centre taps
of the signal transformers so that power and differential data coexist. After a detection handshake
(the powered device presents a 25\,k$\Omega$ signature) and classification, the switch applies about
44 to 57\,V. IEEE 802.3af (2003) provided up to 15.4\,W at the switch port, 802.3at (2009) 30\,W, and
802.3bt (2018) 60 and 90\,W using all four pairs, enough for pan-tilt cameras, Wi-Fi access points,
small-cell radios and LED lighting. A line-powered device is another echo of the telephone network,
whose loops have carried $-48$\,V to telephones since the 1890s.

%======================================================================
\section{Optical fibre}
\label{sec:ch24:fibre}
%======================================================================

\begin{history}[From light pipes to the transoceanic fibre]
Guiding light by total internal reflection in a stream of water was demonstrated in the 1840s, and
glass fibres were used for imaging from the 1950s, but their loss, around 1000\,dB/km, made them useless
for communication. In 1966 Charles Kao and George Hockham, at Standard Telecommunication Laboratories in
Harlow, England, argued that this loss came from impurities, mostly iron and other transition metals,
rather than from the glass itself, and that a fibre with less than 20\,dB/km would be a practical
transmission medium. In 1970 Robert Maurer, Donald Keck and Peter Schultz at Corning made fused-silica
fibre that met the target, the same year in which the first semiconductor lasers operated continuously
at room temperature. By 1979 researchers at NTT had reached 0.2\,dB/km at 1550\,nm, close to the
fundamental limit. Commercial systems followed quickly: multimode links at 850\,nm in the late 1970s,
single-mode links at 1310\,nm in the early 1980s. TAT-8, the first transatlantic fibre cable (1988),
carried 280\,Mb/s per fibre pair at 1310\,nm through regenerating repeaters.

Then came the amplifier. In 1987 groups at the University of Southampton (Mears, Reekie, Jauncey and
Payne) and at Bell Labs (Desurvire, Simpson and Becker) demonstrated the erbium-doped fibre amplifier,
which amplifies light directly over a band of some 30\,nm around 1550\,nm. Many wavelengths could now be
amplified together, and wavelength-division multiplexing took off: TAT-12/13 (1996) was the first
transatlantic cable with optical amplifiers, and by 2000 commercial systems carried 80 or more
wavelengths of 10\,Gb/s. The capacity boom of the late 1990s ended in the telecom crash of 2001,
leaving a glut of unused (``dark'') fibre that took a decade to fill. Kao received the Nobel Prize in
Physics in 2009.
\end{history}

\subsection{Guiding light}

An optical fibre is a cylindrical waveguide of fused silica: a \emph{core} of refractive index $n_1$
surrounded by a \emph{cladding} of slightly lower index $n_2$, the difference being made by doping the
core with germanium (or the cladding with fluorine). Light striking the core--cladding boundary at
more than the critical angle is totally reflected, so the fibre accepts rays within a cone whose
half-angle is given by the \term{numerical aperture},
\begin{equation}
\mathrm{NA}=\sin\theta_{\max}=\sqrt{n_1^2-n_2^2}\approx n_1\sqrt{2\Delta},\qquad
\Delta=\frac{n_1-n_2}{n_1}.
\end{equation}
For standard single-mode fibre $\Delta\approx0.35\%$ and $\mathrm{NA}\approx0.12$; for multimode fibre
$\Delta\approx1\%$ and $\mathrm{NA}\approx0.2$.

The ray picture is only an approximation: the fibre supports a discrete set of \emph{modes}, field
patterns that propagate unchanged, and their number is governed by the normalised frequency
\begin{equation}
V=\frac{2\pi a}{\lambda}\,\mathrm{NA},
\end{equation}
where $a$ is the core radius. When $V<2.405$ (the first zero of the Bessel function $J_0$), only the
fundamental mode propagates and the fibre is \term{single-mode}. ITU-T G.652 standard single-mode fibre
(SMF) has a core about 8 to 9\,$\mu$m in diameter, a cable cutoff wavelength below 1260\,nm and a mode
field diameter of about 9\,$\mu$m at 1310\,nm and 10.4\,$\mu$m at 1550\,nm, corresponding to an effective
area of about 80\,$\mu$m$^2$. \term{Multimode} fibre has a 50 or 62.5\,$\mu$m core and carries hundreds of
modes; its large core makes connectors and light sources cheap, but its modes travel at different speeds.

\term{Modal dispersion} is the spread in arrival time between modes. In a step-index multimode fibre the
fastest ray runs along the axis and the slowest at the critical angle, giving a spread of
$\Delta\tau\approx Ln_1\Delta/c$: about 50\,ns per kilometre for $\Delta=1\%$, which limits the rate to
a few tens of megabits per second over a kilometre. A \emph{graded-index} core, with the index falling
parabolically from the axis, makes the longer paths travel through faster glass and reduces the spread
by a factor of about $\Delta/8$. Modern laser-optimised multimode fibres are specified by their
\emph{effective modal bandwidth}: 2000\,MHz$\cdot$km for OM3 and 4700\,MHz$\cdot$km for OM4 at 850\,nm,
enough for 100\,Gb/s-per-lane links over about 100\,m. Inside data centres, multimode fibre with cheap
850\,nm VCSELs remains popular for short links; everything longer is single-mode.

\subsection{Attenuation}

The loss of silica fibre (Figure~\ref{fig:ch24_fibre_loss}) is set by three mechanisms.
\term{Rayleigh scattering} from frozen-in density fluctuations of the glass, on a scale far smaller than
the wavelength, scales as $\lambda^{-4}$ and dominates at short wavelengths. \emph{Infrared absorption}
by the vibrations of the silicon--oxygen bonds rises steeply above about 1600\,nm. Between them lies a
minimum near 1550\,nm, where modern fibre loses 0.17 to 0.2\,dB/km and ultra-low-loss pure-silica-core
fibres approach 0.15\,dB/km. Impurities add absorption peaks, the most important being the hydroxyl
(OH$^-$) overtone at 1383\,nm; ``low-water-peak'' fibre (G.652.D) dries the glass so thoroughly that the
peak almost disappears and the whole range from 1260 to 1625\,nm is usable.

\mdcfig[0.85]{ch24_fibre_loss}{Attenuation of silica single-mode fibre (a simple model fitted to typical
values). Rayleigh scattering ($\propto\lambda^{-4}$) dominates at short wavelengths and infrared absorption
beyond 1600\,nm, leaving a minimum of about 0.18\,dB/km near 1550\,nm. Legacy fibre has a water peak at
1383\,nm, which low-water-peak fibre removes. The shaded regions are the O, E, S, C, L and U bands.}

The usable spectrum is divided into bands: O (original, 1260--1360\,nm), E (extended, 1360--1460\,nm),
S (short, 1460--1530\,nm), C (conventional, 1530--1565\,nm), L (long, 1565--1625\,nm) and U (ultra-long,
1625--1675\,nm). Early single-mode systems used the O band, where dispersion is nearly zero, and short-reach
links still do. Long-haul systems live in the C band, where the loss is lowest and erbium amplifiers work,
increasingly extended to C+L. The full low-loss window of a modern fibre is about 50\,THz wide, far more
than any transmission system yet uses.

\begin{inpractice}[Fibre counts and what a link loses]
A span budget needs more than the fibre's dB/km. Fusion splices lose about 0.02 to 0.1\,dB each, and a
cable is spliced every 2 to 6\,km; connectors lose 0.2 to 0.5\,dB each; and planners add an ageing and
repair margin of a few decibels, because every cable cut is repaired with extra splices. In practice a
terrestrial span is budgeted at about 0.22 to 0.25\,dB/km all-in. Cables themselves carry from a few to
several thousand fibres (ribbon cables with 3456 or 6912 fibres are common between data centres), so fibre
is rarely the scarce resource---the expensive parts are the trench, the amplifiers and the transceivers.
\end{inpractice}

\subsection{Chromatic dispersion}

In single-mode fibre there is only one mode, but its group velocity depends on wavelength, both because
the index of silica does (\emph{material} dispersion) and because the fraction of the mode's power in the
core does (\emph{waveguide} dispersion). The \term{chromatic dispersion} parameter
\begin{equation}
D=\frac{d}{d\lambda}\left(\frac{1}{v_g}\right)=-\frac{2\pi c}{\lambda^2}\beta_2
\label{eq:ch24:D}
\end{equation}
measures the differential delay per unit wavelength and length, in ps/(nm$\cdot$km), and $\beta_2$ is the
group-velocity dispersion of Chapter~\ref{ch:equalization}. A signal of spectral width $\Delta\lambda$ is
spread in time by
\begin{equation}
\Delta\tau=|D|\,L\,\Delta\lambda,\qquad \Delta\lambda\approx\frac{\lambda^2}{c}R_s,
\label{eq:ch24:spread}
\end{equation}
where the second form applies to a modulated signal of symbol rate $R_s$ whose spectrum is wider than the
laser line. In G.652 fibre $D$ passes through zero between 1300 and 1324\,nm and reaches about
17\,ps/(nm$\cdot$km) at 1550\,nm (Figure~\ref{fig:ch24_dispersion}). Because the delay spread grows with
$R_s$ through $\Delta\lambda$ and the symbol period shrinks as $1/R_s$, the spread measured in symbols grows
as $R_s^2$: doubling the symbol rate quarters the distance a given dispersion tolerance allows.

\begin{worked}[Chromatic dispersion over an 80\,km span at 32\,GBd]
A 32\,GBd signal at 1550\,nm crosses one 80\,km span of G.652 fibre. How much does dispersion spread it?

The accumulated dispersion is $DL=17\times80=1360$\,ps/nm. The signal's spectral width is
$\Delta\lambda=\lambda^2R_s/c=(1.55\times10^{-6})^2(32\times10^9)/(3\times10^8)=0.256$\,nm, so
$\Delta\tau=1360\times0.256\approx350$\,ps. The symbol period is 31.25\,ps, so each pulse is smeared over
about 11 symbols: the received constellation is a featureless cloud (Figure~\ref{fig:ch24_coherent_dsp}).
For a 10\,Gb/s NRZ signal the same span spreads pulses by about 108\,ps, about one bit, which is roughly
the limit that intensity-modulated 10\,Gb/s links can tolerate without compensation. A coherent receiver
undoes 1360\,ps/nm with a fixed all-pass filter of a few dozen taps; a transpacific link accumulates
more than 150\,000\,ps/nm, which the same receiver also removes, with a filter thousands of taps long
(Section~\ref{sec:ch12:optical}).
\end{worked}

\mdcfig{ch24_dispersion}{Left: the dispersion parameter of standard (G.652), dispersion-shifted (G.653)
and non-zero dispersion-shifted (G.655, an approximate example) fibres. Right: the small-signal response
of a chirp-free intensity-modulated, directly detected link at 1550\,nm. Because the photodiode detects
power, the two sidebands of the intensity modulation, phase-rotated in opposite directions by
dispersion, interfere and cancel at frequencies $f_k=\sqrt{(2k+1)c/(2DL\lambda^2)}$. After 80\,km the
first null is at 6.8\,GHz, which is why 10\,Gb/s IM/DD links reach about 80\,km and 50--100\,GBd links
are confined to a few kilometres or to the O band.}

\subsubsection{Why direct detection hates dispersion}

For an intensity-modulated signal, the damage is worse than pulse spreading suggests. A small-signal
intensity modulation at frequency $f$ places two field sidebands at $\pm f$ around the carrier. Dispersion
rotates their phases by equal amounts $\phi=\beta_2L(2\pi f)^2/2$ in the same sense, and when the
photodiode squares the field the two beat notes with the carrier combine as $\cos\phi$:
\begin{equation}
H_{\mathrm{IM/DD}}(f)=\cos\!\left(\frac{\pi D L\lambda^2f^2}{c}\right).
\label{eq:ch24:fading}
\end{equation}
The response has deep nulls (Figure~\ref{fig:ch24_dispersion}, right), and because the detector has
discarded the phase, no receiver equaliser can fill them. That is why the 80\,km limit of 10\,Gb/s
links was overcome before 2010 by inserting \emph{dispersion-compensating fibre} (with a large negative
$D$) at every amplifier site, why modern data-centre links at 50 to 100\,GBd run in the O band, and why
long-haul systems moved to coherent detection, which keeps the field and turns dispersion into a
harmless all-pass filter.

\begin{pitfall}[Zero dispersion is not always good]
It seems natural to design fibre with zero dispersion at 1550\,nm, and G.653 dispersion-shifted fibre,
laid widely in the late 1980s and 1990s (especially in Japan and Italy), did exactly that. It was ideal
for single-wavelength 10\,Gb/s systems but turned out to be poor for WDM: with no dispersion, neighbouring
wavelengths stay in phase along the fibre, and the nonlinear \emph{four-wave mixing} between them builds
up coherently, creating interference products on top of other channels. The industry's response was
non-zero dispersion-shifted fibre (G.655), with a small but deliberately non-zero $D$ of a few
ps/(nm$\cdot$km) in the C band. In the coherent era, where the DSP removes any amount of dispersion,
\emph{large} dispersion is actually beneficial, because it decorrelates the channels and weakens all
nonlinear interactions; new long-haul builds use G.652 or large-effective-area G.654 fibre.
\end{pitfall}

\subsection{Polarisation-mode dispersion}

A real fibre is not perfectly circular, and stress and bending make it slightly birefringent, with an
orientation that varies randomly along the fibre and with time. The result is
\term{polarisation-mode dispersion} (PMD): the signal splits between two principal states of polarisation
that arrive with a \emph{differential group delay} (DGD). Because the birefringent sections add like a
random walk, the mean DGD grows as $\sqrt L$: modern fibre has a PMD coefficient below about
0.1\,ps/$\sqrt{\text{km}}$, so a 1000\,km link has a mean DGD of about 3\,ps, while older fibre laid in the
1980s can be ten times worse. The DGD follows a Maxwellian distribution, so a link must be designed for
an outage probability, typically allowing an instantaneous DGD three times the mean. PMD was a serious
limit for 40\,Gb/s direct-detection systems; in coherent receivers it is just part of the $2\times2$
MIMO channel removed by the butterfly equaliser (Section~\ref{sec:ch12:optical}).

\subsection{Nonlinearity}

Silica is nonlinear only very weakly, but in a fibre the light is confined to an area of 80\,$\mu$m$^2$ for
tens of kilometres, and the effects add up. The dominant mechanism is the \term{Kerr effect}, an
intensity-dependent refractive index $n=n_0+n_2I$ with $n_2\approx2.6\times10^{-20}$\,m$^2$/W. The
nonlinear phase acquired over a span is
\begin{equation}
\phi_{\mathrm{NL}}=\gamma P L_{\mathrm{eff}},\qquad
\gamma=\frac{2\pi n_2}{\lambda A_{\mathrm{eff}}}\approx1.3\ \mathrm{W^{-1}km^{-1}},\qquad
L_{\mathrm{eff}}=\frac{1-e^{-\alpha L}}{\alpha}\approx\frac1\alpha\approx22\ \text{km},
\end{equation}
so a channel launched at 1\,mW acquires about 0.03\,rad per span; with a hundred channels and dozens of
spans, these small phases become the main impairment. The Kerr effect appears in three guises.
\emph{Self-phase modulation} (SPM): each channel's own intensity modulates its phase, which dispersion
converts into amplitude distortion. \emph{Cross-phase modulation} (XPM): the intensity of every other
channel modulates the phase of this one. \emph{Four-wave mixing} (FWM): three frequencies $f_i+f_j-f_k$
mix to create a new one, on top of another channel if the grid is regular. Two inelastic scattering
processes also matter. \emph{Stimulated Brillouin scattering} reflects power from a narrow-linewidth
carrier above a threshold of a few milliwatts, which is why analog cable-TV transmitters dither their
lasers. \emph{Stimulated Raman scattering} transfers power from short- to long-wavelength channels
across a wideband C+L system, tilting the spectrum, and is also exploited in Raman amplifiers.

\begin{keyidea}[The nonlinear Shannon limit]
In a linear channel, raising the launch power always raises the SNR. In a fibre it does so only until
the Kerr nonlinearity, which grows as the cube of the power, takes over. A useful model, the Gaussian-noise
(GN) model of Poggiolini and co-workers, treats the nonlinear interference as additional Gaussian noise
of power $\eta P^3$, so the SNR after $N$ spans is
\begin{equation}
\mathrm{SNR}=\frac{P}{N\left(P_{\mathrm{ASE}}+\eta P^3\right)}.
\label{eq:ch24:gn}
\end{equation}
Setting the derivative to zero gives the optimum launch power $P_{\mathrm{opt}}=(P_{\mathrm{ASE}}/2\eta)^{1/3}$,
at which the nonlinear noise is exactly half the amplifier noise, and the best SNR,
$\mathrm{SNR}_{\max}=\tfrac23P_{\mathrm{opt}}/(NP_{\mathrm{ASE}})$, which scales as $(NP_{\mathrm{ASE}})^{-2/3}$:
doubling the distance costs 2\,dB, not 3\,dB, and halving the amplifier noise gains 2\,dB, not 3. The
capacity of a fibre, $2\log_2(1+\mathrm{SNR}_{\max})$ per polarisation pair per symbol, therefore has a
maximum, the so-called \emph{nonlinear Shannon limit} (Figure~\ref{fig:ch24_osnr_nli}, right). It is not a
true capacity bound---the nonlinear interference is deterministic in principle and partly
compensable---but it is the practical ceiling that every modern system design works against.
\end{keyidea}

%======================================================================
\section{Optical transmitters and receivers}
\label{sec:ch24:trx}
%======================================================================

\subsection{Light sources and modulators}

The light source of nearly every fibre system is a semiconductor laser. A \emph{Fabry--Perot} laser, a
gain chip between two cleaved facets, oscillates on several longitudinal modes at once and is used
only for short, low-rate links. The \emph{distributed-feedback} (DFB) laser etches a Bragg grating along
its active region so that only one mode lases; its linewidth is typically a few hundred kilohertz to a
few megahertz. \emph{External-cavity} and widely tunable lasers (the integrable tunable laser assembly,
ITLA, of coherent modules) reach linewidths of 100\,kHz or below, which coherent QAM needs for its carrier
recovery. The \emph{vertical-cavity surface-emitting laser} (VCSEL) emits from its top surface, can be
tested on the wafer and made in arrays, and at 850\,nm is the cheap source of multimode data-centre links.

The simplest transmitter modulates the laser's drive current directly. Directly modulated lasers are
cheap, but modulating the carrier density also modulates the refractive index, so the optical frequency
swings with the intensity (\emph{chirp}), broadening the spectrum and worsening dispersion. Faster and
cleaner transmitters keep the laser running continuously and modulate its output externally. An
\emph{electro-absorption modulator} (EAM), often integrated with the DFB as an EML, changes its absorption
with applied field. The \term{Mach--Zehnder modulator} (MZM) splits the light into two arms whose phases
are shifted by electro-optic effect in lithium niobate, indium phosphide or silicon, and recombines them.
Driven in push--pull with voltage $V$, its field transfer is
\begin{equation}
\frac{E_{\mathrm{out}}}{E_{\mathrm{in}}}=\cos\!\left(\frac{\pi V}{2V_\pi}\right),
\end{equation}
where $V_\pi$ is the voltage for a $\pi$ phase difference. Biased at the half-power point it is a
chirp-free intensity modulator; biased at the null it becomes a \emph{field} modulator whose output goes
through zero and changes sign, producing BPSK or, in pairs, the IQ modulation of coherent transmitters
(Figure~\ref{fig:ch24_coherent}). Thin-film lithium niobate and silicon photonic modulators now reach
analog bandwidths of 50 to more than 100\,GHz.

\subsection{Photodetectors and noise}

A \term{photodiode} converts optical power into current with a \term{responsivity}
\begin{equation}
\mathcal R=\frac{\eta q}{h\nu}=\eta\,\frac{\lambda[\mu\mathrm m]}{1.24}\ \ \text{A/W},
\end{equation}
where $\eta$ is the quantum efficiency; an ideal InGaAs detector gives 1.25\,A/W at 1550\,nm and
practical ones 0.8 to 1.0\,A/W. The current carries three kinds of noise. \emph{Shot noise}, from the
discreteness of photons and electrons, has a one-sided spectral density $2qI$ (Chapter~\ref{ch:noise}).
\emph{Thermal noise} of the load or transimpedance amplifier (TIA) has density $4kT/R_L$, or is specified
directly as an input-referred noise current density of typically 10 to 20\,pA/$\sqrt{\text{Hz}}$ for
multi-gigabit TIAs. And in an amplified system, the \emph{beat noise} between the signal and the
amplified spontaneous emission of optical amplifiers (Section~\ref{sec:ch24:wdm}) dominates both.

For on--off keying, with photocurrents $I_1$ and $I_0$ and noise standard deviations $\sigma_1$ and
$\sigma_0$ for the two levels, the optimum threshold gives
\begin{equation}
\mathrm{BER}=\tfrac12\operatorname{erfc}\!\left(\frac{Q}{\sqrt2}\right),\qquad Q=\frac{I_1-I_0}{\sigma_1+\sigma_0},
\label{eq:ch24:qfactor}
\end{equation}
so that $Q=6$ gives $10^{-9}$ and $Q=7.03$ gives $10^{-12}$; optical engineers often quote $20\log_{10}Q$
in decibels. The \term{receiver sensitivity} is the average power needed for a given BER.

An unamplified \emph{PIN} receiver is limited by thermal noise. An \term{avalanche photodiode} (APD)
multiplies the photocurrent by an internal gain $M$, which lifts the signal above the thermal noise, but
the multiplication is random and multiplies the shot noise by $M^2F(M)$, with excess noise factor
\begin{equation}
F(M)=k_AM+(1-k_A)\left(2-\frac1M\right),
\end{equation}
where $k_A$ is the ratio of hole and electron ionisation rates (about 0.02 for silicon, 0.4 to 0.5 for
InGaAs/InP APDs). Increasing $M$ first raises the SNR, as the signal grows above the fixed thermal noise,
and then lowers it, as excess shot noise grows faster than the signal; there is an optimum gain,
typically 5 to 20 (Figure~\ref{fig:ch24_rx_sensitivity}). APDs gain roughly 6 to 10\,dB of sensitivity
over PINs, which is why they are used in PON optical line terminals and long unamplified links.

The ultimate reference is the \emph{quantum limit}: an ideal photon counter with no thermal noise makes an
error only when a one-bit produces no photons at all, with probability $e^{-N_1}$ for a mean of $N_1$
photons. For $10^{-9}$, $N_1=20$, an average of 10 photons per bit.

\mdcfig{ch24_rx_sensitivity}{Left: bit error rate of 10\,Gb/s on--off keying at 1550\,nm against
received power, for a thermal-noise-limited PIN receiver (input noise 15\,pA/$\sqrt{\text{Hz}}$,
responsivity 0.9\,A/W), an InGaAs APD with $k_A=0.5$ at its optimum gain, and the quantum limit. Right:
the APD's Q-factor against gain at three received powers; the optimum gain balances thermal noise against
excess avalanche noise.}

\begin{worked}[How far from the quantum limit is a real receiver?]
For the 10\,Gb/s receiver of Figure~\ref{fig:ch24_rx_sensitivity}, the PIN version reaches $10^{-12}$ at
about $-20$\,dBm and the APD version at about $-29$\,dBm. How many photons per bit are these?

At 1550\,nm a photon carries $h\nu=hc/\lambda=1.28\times10^{-19}$\,J. At $-20$\,dBm, $10\,\mu$W, the
energy per bit is $10^{-5}/10^{10}=10^{-15}$\,J, or about 7800 photons. The APD needs about 980. The
quantum limit at $10^{-12}$ is about 14 photons per bit ($e^{-N_1}=2\times10^{-12}$ gives $N_1\approx27$,
half of that on average). A PIN receiver is therefore some 27\,dB, and an APD some 18\,dB, from the
quantum limit: the price of thermal noise in the electronics. Optical preamplification with an EDFA
closes much of the gap, to within about 3\,dB of the quantum limit for amplified on--off keying, and
coherent detection with a strong local oscillator is shot-noise limited for the same reason: both lift
the signal far above the electronic noise \emph{before} it is detected.
\end{worked}

\subsection{Intensity-modulated links and their budgets}

An unamplified \emph{intensity-modulation, direct-detection} (IM/DD) link is designed with a simple power
budget: the transmitter's launch power minus the receiver sensitivity must exceed the sum of fibre loss,
connector and splice losses, \emph{power penalties} for impairments (dispersion, finite extinction ratio,
reflections, crosstalk) and a margin for ageing and repairs. Standards such as IEEE 802.3 specify each
term so that any compliant transmitter works with any compliant receiver over the specified fibre. Short
links are dispersion-limited rather than loss-limited at high rates, and the budget therefore always
includes a dispersion penalty computed for the worst-case wavelength. The worked example of
Section~\ref{sec:ch24:pon} applies such a budget to a passive optical network.

%======================================================================
\section{Optical amplifiers and wavelength-division multiplexing}
\label{sec:ch24:wdm}
%======================================================================

\subsection{The erbium-doped fibre amplifier}

An \term{erbium-doped fibre amplifier} (EDFA) is a few to a few tens of metres of fibre whose core is
doped with erbium ions, pumped by a semiconductor laser at 980 or 1480\,nm. The pump raises the ions to an
excited level; the signal photons, in the 1530--1565\,nm band, stimulate their return to the ground state,
each emitting a photon identical to the stimulating one. Gains of 20 to 30\,dB with output powers of
20\,dBm or more are routine, the gain is almost independent of the modulation (the excited state lives
about 10\,ms, so the amplifier responds only to the average power) and all wavelengths in the band are
amplified together. L-band EDFAs use longer, differently designed doped fibre to cover 1565--1625\,nm.
\emph{Raman amplifiers} instead pump the transmission fibre itself about 100\,nm below the signal and
provide distributed gain along the span, which improves the noise performance because the signal is
amplified before it has decayed to its lowest level.

Amplification has an unavoidable cost: excited ions also decay \emph{spontaneously}, emitting photons in
random directions and phases, and those emitted into the guided mode are amplified along with the signal.
This \term{amplified spontaneous emission} (ASE) is the noise of optical systems. Its power spectral
density in each polarisation at the amplifier output is
\begin{equation}
S_{\mathrm{ASE}}=n_{\mathrm{sp}}h\nu(G-1),
\end{equation}
where $G$ is the gain and $n_{\mathrm{sp}}\ge1$ is the spontaneous emission factor, equal to one only with
complete population inversion. The amplifier's \term{noise figure}, defined as for electronic amplifiers
by the degradation of the shot-noise-limited SNR, is $\mathrm{NF}\approx2n_{\mathrm{sp}}$ for large gain:
no phase-insensitive optical amplifier can have a noise figure below 3\,dB, a quantum limit related to the
uncertainty principle. Practical EDFAs achieve 4.5 to 6\,dB.

\subsection{The OSNR budget}

\begin{figure}[htbp]\centering
\begin{tikzpicture}[font=\sffamily\scriptsize,
  box/.style={draw=navy,fill=navy!6,thick,rounded corners=2pt,minimum height=0.7cm,minimum width=0.9cm,align=center},
  amp/.style={draw=accent,fill=accent!8,thick,isosceles triangle,isosceles triangle apex angle=60,inner sep=0.5pt,minimum size=0.55cm},
  fib/.style={thick,accent},
  arr/.style={-{Stealth[length=1.6mm]},thick,navy}]
\foreach \i/\y in {1/0.75,2/0.25,3/-0.25,4/-0.75} {\node[box,minimum height=0.38cm,minimum width=0.75cm] (tx\i) at (0,\y) {Tx $\lambda_\i$};}
\node[box,minimum height=1.9cm,minimum width=0.6cm] (mux) at (1.25,0) {M\\U\\X};
\foreach \i in {1,...,4} \draw[arr] (tx\i.east) -- (mux.west |- tx\i);
\node[amp] (boost) at (2.25,0) {};
\node[below=0.25cm of boost] {booster};
\draw[fib] (mux.east) -- (boost.west);
\draw[fib] (boost.east) -- (2.9,0);
\draw[fib] (3.25,0) circle (0.35); \node at (3.25,-0.6) {80\,km};
\draw[fib] (3.6,0) -- (4.0,0);
\node[amp] (in1) at (4.3,0) {}; \node[below=0.25cm of in1] {in-line};
\draw[fib] (in1.east) -- (4.95,0);
\draw[fib] (5.3,0) circle (0.35); \node at (5.3,-0.6) {80\,km};
\draw[fib] (5.65,0) -- (6.0,0);
\node[box,fill=accent!6,draw=accent,minimum height=1.4cm,minimum width=1.3cm] (roadm) at (6.75,0) {ROADM\\(WSS)};
\draw[arr] (6.5,-1.25) -- node[left]{add} (6.5,-0.72); \draw[arr] (7.0,-0.72) -- node[right]{drop} (7.0,-1.25);
\draw[fib] (7.4,0) -- (7.8,0);
\node at (8.15,0) {$\cdots$};
\node[amp] (pre) at (8.75,0) {}; \node[below=0.25cm of pre] {pre-amp};
\draw[fib] (pre.east) -- (9.45,0);
\node[box,minimum height=1.9cm,minimum width=0.6cm] (dmx) at (9.75,0) {D\\E\\M\\U\\X};
\foreach \i/\y in {1/0.75,2/0.25,3/-0.25,4/-0.75} {\node[box,minimum height=0.38cm,minimum width=0.75cm] (rx\i) at (11.05,\y) {Rx $\lambda_\i$}; \draw[arr] (dmx.east |- rx\i) -- (rx\i.west);}
\draw[decorate,decoration={brace,amplitude=4pt},navy] (2.0,1.05) -- node[above=4pt]{$N$ spans, each with loss $L_s\approx16$--20\,dB and an amplifier of gain $G=L_s$} (8.9,1.05);
\end{tikzpicture}
\caption{A wavelength-division-multiplexed line system. Transmitters on different wavelengths are
multiplexed onto one fibre and amplified together by a booster, in-line amplifiers every 50 to 100\,km
and a preamplifier. Reconfigurable optical add--drop multiplexers (ROADMs), built from wavelength-selective
switches, add, drop or pass through individual wavelengths without electrical conversion.}
\label{fig:ch24_wdm}
\end{figure}

A long-haul line (Figure~\ref{fig:ch24_wdm}) consists of $N$ spans, each with loss $L_s$ (in linear units)
followed by an amplifier of gain $G=L_s$ that restores the power. Each amplifier adds ASE of power
$\mathrm{NF}\cdot h\nu G B$ in a bandwidth $B$ (both polarisations, high gain), and the contributions of
identical spans add. The figure of merit is the \term{optical signal-to-noise ratio} (OSNR), conventionally
the ratio of the channel power to the ASE power in a reference bandwidth of 0.1\,nm, or 12.5\,GHz at
1550\,nm:
\begin{equation}
\mathrm{OSNR}_{\mathrm{dB}}=P_{\mathrm{ch}}[\mathrm{dBm}]-L_s[\mathrm{dB}]-\mathrm{NF}[\mathrm{dB}]
-10\log_{10}N+58,
\label{eq:ch24:osnr}
\end{equation}
where $58=-10\log_{10}(h\nu B_{\mathrm{ref}}/1\,\mathrm{mW})$ at 1550\,nm. Each decibel of span loss or noise
figure costs a decibel of OSNR, and each doubling of the number of spans costs 3\,dB. The OSNR is converted
to the electrical SNR per symbol that the modulation formats of Chapter~\ref{ch:modulation} need by
$\mathrm{SNR}=\mathrm{OSNR}\cdot B_{\mathrm{ref}}/R_s$ for a dual-polarisation signal of symbol rate $R_s$.

\begin{worked}[An OSNR budget for a 1000\,km link]
A 1000\,km terrestrial link uses 80\,km spans budgeted at 0.22\,dB/km, so $L_s=17.6$\,dB, with EDFAs of 5\,dB
noise figure. Thirteen amplified spans (the last slightly short) cover the distance. The launch power is
1\,dBm per channel. Can it carry 400\,Gb/s as dual-polarisation 16-QAM at 64\,GBd?

From~\eqref{eq:ch24:osnr}: $\mathrm{OSNR}=1-17.6-5-10\log_{10}13+58=1-17.6-5-11.1+58=25.3$\,dB in 0.1\,nm.
Converting to the 64\,GBd signal bandwidth subtracts $10\log_{10}(64/12.5)=7.1$\,dB, so the ASE-limited SNR
is 18.2\,dB. Near the optimum launch power the nonlinear interference adds half as much noise again
(the key idea above), a further $10\log_{10}1.5=1.8$\,dB, leaving about 16.4\,dB. Dual-polarisation 16-QAM
with a soft-decision FEC that corrects a pre-FEC BER of $2\times10^{-2}$ needs an SNR of about 12.7\,dB in
theory, and about 14.7\,dB allowing a typical 2\,dB of transceiver implementation penalty. The margin is
therefore under 2\,dB, which is tight but workable: 400\,Gb/s at 64\,GBd is a format for links of up to
about 1000\,km. The same link at 32\,GBd with dual-polarisation QPSK (100\,Gb/s), needing about 6.3\,dB of
SNR, has over 10\,dB of margin and could reach several thousand kilometres
(Figure~\ref{fig:ch24_osnr_nli}).
\end{worked}

\mdcfig{ch24_osnr_nli}{Left: linear OSNR from~\eqref{eq:ch24:osnr} against distance for 80\,km spans
(17.6\,dB loss, 5\,dB NF) at three launch powers, with the approximate OSNR required by 100\,G DP-QPSK at
32\,GBd and 400\,G DP-16QAM at 64\,GBd (pre-FEC BER $2\times10^{-2}$ plus 2\,dB implementation margin).
Centre: SNR of a 64\,GBd channel in a fully loaded 4.8\,THz C band against launch power (GN model,
$\gamma=1.3$\,W$^{-1}$km$^{-1}$, $D=17$\,ps/(nm$\cdot$km)); dotted lines omit nonlinearity. The optimum
power is about 1\,dBm per channel. Right: the resulting spectral efficiency limit
$2\log_2(1+\mathrm{SNR}_{\max})$ per 75\,GHz slot against distance.}

\subsection{WDM grids and ROADMs}

Wavelength-division multiplexing puts many independent channels on one fibre, each on its own carrier.
The ITU-T G.694.1 grid anchors \term{dense WDM} (DWDM) channels at 193.1\,THz (about 1552.5\,nm) with
fixed spacings of 100, 50, 25 or 12.5\,GHz; the C band holds about 40 channels at 100\,GHz, or 80 to 96 at
50\,GHz. As symbol rates rose beyond what fits in 50\,GHz, the \emph{flexible grid} of the same
Recommendation allowed any channel centred on a 6.25\,GHz raster with a width that is a multiple of
12.5\,GHz, so that a 64\,GBd channel occupies 75\,GHz and a 130\,GBd channel 150\,GHz. The low-cost
\emph{coarse} WDM grid of G.694.2 uses 18 wavelengths 20\,nm apart from 1271 to 1611\,nm, wide enough for
uncooled lasers, and is used in access and data-centre links that need no amplifiers.

A modern network is not a set of point-to-point lines but a mesh of nodes, each built around a
\term{reconfigurable optical add--drop multiplexer} (ROADM). Its core is the \emph{wavelength-selective
switch} (WSS), in which a diffraction grating spreads the spectrum across a liquid-crystal-on-silicon or
MEMS array that steers each slice of spectrum to any of several output fibres. A ROADM can drop a
wavelength to a local transceiver, add one, or pass it through to another direction optically, without
electrical regeneration, under software control. Modern \emph{colourless, directionless, contentionless}
ROADMs let any transceiver use any wavelength in any direction, which makes the optical layer as flexible
as the packet layer above it. Light paths of several thousand kilometres across many ROADMs are common,
and the OSNR and filtering penalties they accumulate are tracked by the network's planning tools in
exactly the way of the worked example.

%======================================================================
\section{Coherent optical communications}
\label{sec:ch24:coherent}
%======================================================================

\begin{history}[The coherent revolution, twice]
Coherent detection of light, mixing the signal with a local laser exactly as a superheterodyne radio
mixes with a local oscillator (Chapter~\ref{ch:sdr}), was researched intensively in the 1980s because it
promised 10 to 20\,dB better sensitivity than direct detection and channel selection by electrical
filtering. It needed lasers with stable, narrow linewidths and polarisation tracking with optical
hardware, and it was abandoned almost overnight when the EDFA arrived around 1990: an optical
preamplifier gave direct detection most of the sensitivity benefit at a fraction of the complexity.
For fifteen years, fibre systems carried on--off keying, with dispersion and PMD managed optically.

The revival came from CMOS. By the mid-2000s analog-to-digital converters running at tens of gigasamples
per second and DSP chips of millions of gates had become possible, and several groups, notably Kazuro
Kikuchi's at the University of Tokyo, showed that a \emph{digital} coherent receiver---a free-running local
laser, four ADCs and DSP for phase and polarisation tracking---could replace all the optical tricks of
the 1980s. Nortel shipped a 40\,Gb/s dual-polarisation QPSK coherent system in 2008, and 100\,Gb/s
DP-QPSK at about 28 to 32\,GBd, standardised in an OIF framework, became the industry's long-haul
workhorse from 2010. Dispersion compensation modules disappeared from new networks. The DSP-based
coherent transceiver, now shrunk into a pluggable module the size of a pack of chewing gum, is perhaps
the most complete realisation of this book's ideas in a single product.
\end{history}

\subsection{Transmitter and receiver}

\begin{figure}[htbp]\centering
\begin{tikzpicture}[font=\sffamily\scriptsize,
  box/.style={draw=navy,fill=navy!6,thick,rounded corners=2pt,minimum height=0.6cm,minimum width=1.0cm,align=center},
  opt/.style={draw=accent,fill=accent!6,thick,rounded corners=2pt,minimum height=0.6cm,align=center},
  ol/.style={thick,accent},
  el/.style={-{Stealth[length=1.5mm]},thick,navy}]
% ---------------- transmitter
\node[opt] (las) at (0,2.6) {laser};
\node[opt] (iqx) at (2.6,3.2) {IQ modulator X\\(nested MZMs)};
\node[opt] (iqy) at (2.6,2.0) {IQ modulator Y\\(nested MZMs)};
\node[opt] (pbc) at (5.0,2.6) {PBC};
\draw[ol] (las.east) -- (0.9,2.6) |- (iqx.west); \draw[ol] (0.9,2.6) |- (iqy.west);
\draw[ol] (iqx.east) -- (4.1,3.2) |- (pbc.west); \draw[ol] (iqy.east) -- (4.1,2.0) |- (pbc.west);
\node[box] (dac) at (2.6,4.35) {4 DACs};
\node[box] (txd) at (5.6,4.35) {TX DSP: FEC, mapping,\\shaping, pre-emphasis};
\draw[el] (txd) -- (dac); \draw[el] (dac.south) -- (iqx.north);
\draw[el] (dac.south east) .. controls (4.0,3.8) and (4.0,2.5) .. (iqy.north east);
\draw[ol] (pbc.east) -- (7.0,2.6) node[right,text=accent]{fibre $\rightarrow$};
\node[text=navy,font=\sffamily\scriptsize\bfseries] at (0.3,4.35) {Transmitter};
% ---------------- receiver
\node[opt] (pbs) at (0,0) {PBS};
\node[opt] (lo) at (0,-1.6) {local\\laser (LO)};
\node[opt] (hx) at (2.2,0.5) {90$^\circ$\\hybrid};
\node[opt] (hy) at (2.2,-1.1) {90$^\circ$\\hybrid};
\draw[ol] (-1.2,0) node[left,text=accent]{$\rightarrow$} -- (pbs.west);
\draw[ol] (pbs.east) -- (1.0,0) |- ([yshift=0.12cm]hx.west); \draw[ol] (1.0,0) |- ([yshift=0.12cm]hy.west);
\draw[ol] (lo.east) -- (1.3,-1.6) |- ([yshift=-0.12cm]hx.west); \draw[ol] (1.3,-1.6) |- ([yshift=-0.12cm]hy.west);
\node[box,minimum width=1.4cm] (pdx) at (4.2,0.5) {2 balanced\\PD pairs};
\node[box,minimum width=1.4cm] (pdy) at (4.2,-1.1) {2 balanced\\PD pairs};
\draw[ol] (hx) -- node[above]{4} (pdx); \draw[ol] (hy) -- node[above]{4} (pdy);
\node[box] (adc) at (6.1,-0.3) {4 ADCs\\$\sim$1.2--2\\samples/sym};
\draw[el] (pdx.east) -- (adc.north west); \draw[el] (pdy.east) -- (adc.south west);
\node[box,minimum width=3.4cm,align=left,anchor=west] (dsp) at (7.3,-0.3) {front-end correction\\CD compensation (FFT)\\timing recovery\\$2\times2$ butterfly (CMA, LMS)\\frequency \& phase recovery\\soft demapping, SD-FEC};
\draw[el] (adc) -- (dsp);
\node[text=navy,font=\sffamily\scriptsize\bfseries] at (0.3,1.05) {Receiver};
\end{tikzpicture}
\caption{A dual-polarisation coherent transceiver. The transmitter splits a laser between two IQ
modulators, each a pair of nested Mach--Zehnder modulators driven by DACs, and combines their outputs on
orthogonal polarisations. The receiver splits the incoming light into two polarisations, mixes each with a
local laser in a $90^\circ$ hybrid, detects in-phase and quadrature beat signals with balanced
photodiodes, digitises the four signals and recovers the data in DSP.}
\label{fig:ch24_coherent}
\end{figure}

A coherent transmitter (Figure~\ref{fig:ch24_coherent}) imposes a complex baseband signal on the optical
field. An \emph{IQ modulator} nests two Mach--Zehnder modulators, biased at null, inside a larger
interferometer with a $90^\circ$ phase shift between its arms, so that the output field is proportional to
$I(t)+jQ(t)$ for small drive signals, exactly the quadrature modulator of a radio. Two IQ modulators fed by
one laser and combined on orthogonal polarisations by a polarisation beam combiner give four
dimensions per symbol: dual-polarisation QPSK carries 4 bits per symbol, DP-16QAM 8 and DP-64QAM 12.
Four high-speed DACs, driven by a transmit DSP that maps, shapes (root-raised-cosine with roll-off of a
few per cent) and pre-emphasises the signal, drive the modulators.

The receiver mixes the incoming field $E_s$ with a local oscillator laser $E_{\mathrm{LO}}$ in a
$90^\circ$ \emph{optical hybrid}, a passive interferometer with four outputs $\tfrac12(E_s\pm E_{\mathrm{LO}})$
and $\tfrac12(E_s\pm jE_{\mathrm{LO}})$. A photodiode detects power, $|E_s|^2+|E_{\mathrm{LO}}|^2+2\re\{E_sE_{\mathrm{LO}}^*\}$
up to factors, and subtracting the outputs of two photodiodes (\emph{balanced detection}) cancels the
direct-detection terms and leaves
\begin{equation}
i_I\propto\mathcal R\,\re\{E_sE_{\mathrm{LO}}^*\},\qquad i_Q\propto\mathcal R\,\im\{E_sE_{\mathrm{LO}}^*\}.
\end{equation}
The pair is the complex envelope of the signal, multiplied by the LO amplitude and rotated by the
difference between the two lasers' phases. Three properties follow. The receiver is \emph{linear in the
field}, so dispersion, PMD and polarisation rotation remain linear operations that DSP can invert, and
no information is lost in the detector. The LO acts as a gain of $|E_{\mathrm{LO}}|$ before any electronic
noise, so a strong LO makes the receiver shot-noise limited. And the LO selects the channel: only the
wavelength within the receiver's electrical bandwidth of the LO is detected, so a coherent receiver can
tune to any channel of a WDM comb without an optical filter. Because the LO is a free-running laser
within a few gigahertz of the signal, this is an \emph{intradyne} receiver: the residual frequency and
phase offsets are removed digitally.

\subsection{The DSP chain}

The receive DSP is a tour of Chapters~\ref{ch:sync} and~\ref{ch:equalization}, running on four streams of
samples, each at up to 150 to 200\,GS/s, processed in hundreds of parallel lanes.
\begin{enumerate}
\item \emph{Front-end correction}: deskew between the four channels, compensation of quadrature and gain
imbalance between I and Q (for example by Gram--Schmidt orthogonalisation) and of the frequency response of
the ADCs and photodiodes.
\item \emph{Chromatic dispersion compensation}: a fixed all-pass filter $\exp(+j\beta_2L\omega^2/2)$, applied in
the frequency domain with overlap--save FFTs, its length growing with distance and with the square of the
symbol rate (Section~\ref{sec:ch12:optical}). The accumulated dispersion is estimated blindly at start-up
by sweeping trial values.
\item \emph{Timing recovery}: a Gardner or Godard timing-error detector (Chapter~\ref{ch:sync}) driving an
interpolator.
\item \emph{Polarisation demultiplexing and adaptive equalisation}: the $2\times2$ butterfly of complex FIR
filters of Section~\ref{sec:ch12:optical}, adapted blindly with the constant-modulus algorithm and then by
decision-directed LMS, which undoes polarisation rotation, PMD, residual dispersion and filtering by ROADMs.
\item \emph{Carrier recovery}: estimation of the frequency offset between the signal laser and the LO (for
QPSK, from the peak of the spectrum of the signal raised to the fourth power) and tracking of the laser
phase noise, by the Viterbi--Viterbi fourth-power algorithm for QPSK, by blind phase search for denser
QAM, or with the help of pilot symbols inserted every few dozen symbols.
\item \emph{Soft demapping and decoding}: log-likelihood ratios for each bit feed a soft-decision FEC
(Chapter~\ref{ch:moderncodes}) with 15 to 27\% overhead that corrects pre-FEC bit error rates of
$10^{-2}$ or more to below $10^{-15}$.
\end{enumerate}

\mdcfig{ch24_coherent_dsp}{The coherent DSP chain at work (simulation of 32\,GBd dual-polarisation QPSK,
X polarisation shown): after 80\,km of fibre (1360\,ps/nm), a random polarisation rotation with 8\,ps of
differential group delay, a 150\,MHz laser frequency offset, 200\,kHz combined laser linewidth and
16\,dB SNR. The received samples are a cloud; CD compensation alone leaves the two polarisations mixed;
the blind CMA butterfly restores a constant-modulus ring; frequency-offset removal and Viterbi--Viterbi
phase recovery yield the QPSK constellation. The code is in \lab{commlib/wireline.py}.}

Figure~\ref{fig:ch24_coherent_dsp} shows the chain acting on a simulated 100\,G-class signal. Each stage
removes one impairment, and each does it with an algorithm met earlier in the book; what makes a commercial
receiver hard is doing it at 100 gigabaud with a power budget of a few watts. A coherent DSP ASIC for an
800\,Gb/s module in a 3\,nm or 5\,nm CMOS process performs on the order of $10^{14}$ operations per
second.

\subsection{Pushing towards the limit}

Since 2010 coherent systems have improved along every axis available to them
(Table~\ref{tab:ch24:coherent}). The symbol rate has quadrupled, from about 32 to above 130\,GBd, which
cuts the cost and power per bit roughly in proportion. Constellations have become denser where the SNR
allows, and \term{probabilistic constellation shaping} (PCS) of 64-QAM, transmitting the inner points more
often than the outer ones with a Maxwell--Boltzmann distribution produced by a distribution matcher in the
probabilistic amplitude shaping architecture of B\"ocherer and colleagues, gains up to about 1\,dB of the
1.53\,dB shaping limit (Chapter~\ref{ch:infotheory}) and, more importantly in practice, lets the line rate be
tuned in fine steps to match the SNR of each path. Nokia's PSE-3 chip (2018) brought PCS to commercial
products, and it is now standard in long-haul transceivers.

\begin{table}[htbp]\centering\small
\caption{Generations of coherent transceivers (approximate; vendors and years vary).}
\label{tab:ch24:coherent}
\begin{tabularx}{\linewidth}{@{}l l l X@{}}
\toprule
Rate/$\lambda$ & First products & Format and symbol rate & Notes\\
\midrule
40--100\,Gb/s & 2008--2010 & DP-QPSK, 32\,GBd & First digital coherent generation; long haul\\
200\,Gb/s & 2013 & DP-16QAM or 8QAM, 32\,GBd & Metro and regional\\
400\,Gb/s & 2017--2020 & DP-16QAM, 60--64\,GBd & Includes OIF 400ZR pluggables (2020)\\
800\,Gb/s & 2020--2024 & PCS-64QAM, 90--120\,GBd & 800ZR and 800G long haul\\
1.2--1.6\,Tb/s & 2023--2025 & PCS, 130--200\,GBd & Very high symbol rates, 3\,nm DSP\\
\bottomrule
\end{tabularx}
\end{table}

\begin{inpractice}[400ZR: coherent in a pluggable]
The OIF 400ZR implementation agreement (2020) defined an interoperable coherent interface for data-centre
interconnect: 400\,Gb/s as DP-16QAM at 59.84\,GBd, with a concatenated FEC (an inner soft-decision code
and an outer staircase code, about 15\% overhead; Chapter~\ref{ch:moderncodes}), over amplified links up
to about 120\,km. The whole transceiver, laser, silicon-photonic modulator and receiver, and a 7\,nm DSP,
fits in a QSFP-DD or OSFP module plugged directly into a router, dissipating roughly 15 to 20\,W. Its
successors 800ZR and the OpenZR+ and OpenROADM variants, which use a stronger open FEC (oFEC) for longer
reaches, did to metro transport what VCSEL modules did to the data centre: they made the optical line a
commodity that the routing layer simply plugs in. The pluggable's power budget, not the fibre, now sets
what is possible: a 15\,W module cannot afford the elaborate DSP of a 50\,W line card.
\end{inpractice}

\subsection{Submarine systems}

The longest and most demanding links are the submarine cables, which carry nearly all intercontinental
traffic. A modern transoceanic cable carries 8 to 24 fibre pairs through repeaters spaced every 50 to
80\,km, each containing an EDFA per fibre, powered by a constant direct current fed from both shores
along the cable's copper conductor. The available electrical power, set by the voltage the cable can
withstand, is the binding constraint. Shannon's formula explains what follows: since capacity grows only
logarithmically with SNR but linearly with the number of parallel channels, it is better to run many fibre
pairs at modest power and spectral efficiency than a few pairs at high power. This \term{space-division
multiplexing} (SDM) design philosophy, together with large-effective-area fibre and C-band coherent
transceivers, has taken cable design capacities from a few terabits per second around 2010 to roughly 160
to 250\,Tb/s for cables built around 2017--2021, such as MAREA and Dunant. Laboratory experiments with
multi-core and few-mode fibres have demonstrated petabit-per-second capacities in a single strand,
pointing to the next stage of SDM.

\begin{keyidea}[Capacity per watt, not per hertz]
Radio systems are spectrum-limited and chase bits per hertz. Long-haul fibre systems have tens of terahertz
available and are limited by nonlinearity and by electrical power, so they chase bits per joule. When the
medium is cheap and wide, as in a cable with many fibres, the optimum is to spread the traffic over many
parallel low-SNR channels, the high-SNR wing of the capacity curve being the most expensive place to
operate. The same logic explains why wireless systems add antennas (MIMO, Chapter~\ref{ch:mimo}) rather than
power.
\end{keyidea}

%======================================================================
\section{Passive optical networks}
\label{sec:ch24:pon}
%======================================================================

Running a separate fibre and transceiver from the exchange to every home would be expensive. A
\term{passive optical network} (PON) shares one fibre and one transceiver at the exchange, the
\emph{optical line terminal} (OLT), among 32 to 128 homes by means of unpowered optical splitters in the
field (Figure~\ref{fig:ch24_fttx}). Each home has an \emph{optical network unit} (ONU, or terminal, ONT).
The network is point-to-multipoint, and its two directions behave very differently.
\begin{description}[style=nextline,leftmargin=1.2em,font=\sffamily\bfseries\small\color{navy}]
\item[Downstream: broadcast.] The OLT transmits a continuous TDM stream that the splitter delivers to every
ONU; each ONU keeps the frames addressed to it. Because every home receives every frame, the payload is
encrypted (AES in the ITU-T standards).
\item[Upstream: TDMA.] The ONUs share the fibre towards the OLT in time. The OLT allocates upstream time slots
with a \emph{bandwidth map} in each downstream frame, using \emph{dynamic bandwidth allocation} driven by the
ONUs' reported queue occupancies. Because the ONUs are at different distances---up to 20\,km apart in
fibre---the OLT \emph{ranges} each one at activation, measuring its round-trip delay and assigning an
equalisation delay so that all appear to be at the same distance, the optical counterpart of cable modem
ranging and cellular timing advance. The OLT's receiver is a \emph{burst-mode} receiver: bursts from
different ONUs arrive with power levels differing by up to 15 or 20\,dB and with independent phases, and it
must set its threshold and recover the clock within a preamble of tens of nanoseconds.
\end{description}
The two directions use different wavelengths on the same fibre, separated by WDM filters, and new
generations use new wavelengths so that they can coexist with the old on the same splitter
(Table~\ref{tab:ch24:pon}).

\begin{table}[htbp]\centering\small
\caption{Passive optical network standards (nominal line rates rounded, e.g.\ 2.5 for 2.488\,Gb/s and 10 for 9.953\,Gb/s; years approximate).}
\label{tab:ch24:pon}
\begin{tabularx}{\linewidth}{@{}l l l l X@{}}
\toprule
Standard & Year & Down/up (Gb/s) & $\lambda$ (nm) & Notes\\
\midrule
GPON (G.984) & 2003 & 2.5 / 1.25 & 1490 / 1310 & Mass-market FTTH of the 2010s\\
EPON (802.3ah) & 2004 & 1 / 1 & 1490 / 1310 & Ethernet framing; widely used in Asia\\
10G-EPON (802.3av) & 2009 & 10 / 1 or 10 & 1577 / 1270 & \\
XG-PON (G.987) & 2010 & 10 / 2.5 & 1577 / 1270 & Coexists with GPON on one splitter\\
NG-PON2 (G.989) & 2015 & $4\times10$ / $4\times10$ & 1524--1625 & Tunable TWDM; limited deployment\\
XGS-PON (G.9807.1) & 2016 & 10 / 10 & 1577 / 1270 & Symmetric 10G; the current workhorse\\
50G-PON (G.9804) & 2021 & 50 / 12.5--50 & O band & LDPC FEC, DSP equalisation\\
\bottomrule
\end{tabularx}
\end{table}

The design of a PON is a power-budget exercise, and the standards define \emph{optical budget classes}
that bundle a transmitter power range and a receiver sensitivity: GPON class B+ supports 13 to 28\,dB of
loss between OLT and ONU and class C+ 17 to 32\,dB; XGS-PON classes N1, N2, E1 and E2 support maximum
losses of 29, 31, 33 and 35\,dB. The splitter dominates: an ideal $1{:}N$ split loses $10\log_{10}N$\,dB
(15\,dB for 1:32, 18\,dB for 1:64), and real splitters add 1 to 3\,dB of excess loss.

\begin{worked}[A GPON power budget]
A GPON with class B+ optics serves 32 homes at up to 20\,km. The OLT launches at least $+1.5$\,dBm at
1490\,nm and the ONU's sensitivity is $-27$\,dBm. Will it close downstream?

Losses: fibre $20\times0.30=6.0$\,dB at 1490\,nm; 1:32 splitter about 17.0\,dB including excess loss; four
connectors at 0.3\,dB, 1.2\,dB; ten splices at 0.1\,dB, 1.0\,dB. The total is 25.2\,dB, within the class B+
maximum of 28\,dB. With the minimum launch power, the received power is $1.5-25.2=-23.7$\,dBm, a margin of
3.3\,dB over the sensitivity. Upstream at 1310\,nm the fibre loses about 0.35\,dB/km, adding 1\,dB, and the
OLT's APD receiver (sensitivity about $-28$\,dBm) sees an ONU launching at least $+0.5$\,dBm; the margin is
about 2.3\,dB. Raising the split to 1:64 adds about 3\,dB and exhausts the margin, which is why
high-split or long-reach designs need class C+ optics. Notice how the splitter, not the fibre, dominates:
the fibre loss of a PON is almost an afterthought.
\end{worked}

Today's mass-market fibre-to-the-home is XGS-PON, delivering symmetric 10\,Gb/s per splitter, and 50G-PON
(ITU-T G.9804.3) is entering service. At 50\,Gb/s the PON inherits the techniques of data-centre optics:
it operates in the O band to limit dispersion, needs DSP equalisation in the ONU to cope with the bandwidth
of low-cost optics, and uses an LDPC code to tolerate a pre-FEC BER of about $10^{-2}$.

%======================================================================
\section{Data-centre optics and free-space optics}
\label{sec:ch24:dc}
%======================================================================

The largest volume of optical transceivers by far is inside data centres, where hundreds of thousands of
servers are connected through several tiers of switches by links of a few metres to a few kilometres. The
requirements are brutal: lowest cost and power per bit, enormous volume, and a refresh every two to three
years as switch chips double in capacity (51.2\,Tb/s around 2023, 102.4\,Tb/s around 2025).

\subsection{PAM-4 intensity modulation}

Intra-data-centre links use IM/DD with the PAM-4 signalling of Chapter~\ref{ch:baseband}. A 100\,Gb/s lane
runs at 53.125\,GBd; 400GBASE-DR4 sends four such lanes over four parallel single-mode fibre pairs up to
500\,m, and 400GBASE-FR4 multiplexes four lanes on a 20\,nm CWDM grid in the O band over one fibre pair to
2\,km. The next generation, standardised in IEEE 802.3dj, uses 200\,Gb/s lanes at about 113\,GBd for 800G
and 1.6T ports. The O band is chosen because the power fading of~\eqref{eq:ch24:fading} would close the eye
of a 53 to 113\,GBd intensity-modulated signal after a few kilometres in the C band
(Figure~\ref{fig:ch24_dispersion}); near 1310\,nm the dispersion is almost zero. Each module contains a DSP
that re-times and equalises the electrical signal from the host and the optical signal from the fibre (an
FFE of tens of taps and sometimes a short sequence detector, as in Chapter~\ref{ch:equalization}) with a
Reed--Solomon or concatenated FEC.

\subsection{Linear drive and co-packaged optics}

At 800\,Gb/s per module, the module DSP consumes a large part of the 12 to 16\,W module budget, and a switch
with 64 such ports spends kilowatts on optics. Two architectures attack this. \emph{Linear pluggable optics}
(LPO) remove the DSP from the module altogether: the switch chip's own SerDes, with its strong equaliser,
drives the modulator directly through linear amplifiers, and receives directly from the photodiode's
linear TIA, roughly halving the module power at the cost of a harder link-budget negotiation between host
and module. \emph{Co-packaged optics} (CPO) go further, placing optical engines on the same package substrate
as the switch ASIC, so that the electrical link between them is millimetres rather than tens of centimetres
long and needs little equalisation; the lasers, which dislike heat, sit in separate replaceable external
modules. The trend is the same one that ran through the copper story: shorten the electrical channel, and
the signal processing needed to overcome it shrinks too.

\subsection{Free-space optics}

Light need not be guided. \emph{Free-space optical} (FSO) links send a laser beam through the air between
buildings or across rivers, offering fibre-like rates without a trench and without spectrum licensing. Their
nemesis is weather: fog droplets are comparable in size to the wavelength and attenuate the beam by tens to
hundreds of decibels per kilometre, so terrestrial FSO is limited to short ranges or must be backed up by
radio, while atmospheric turbulence causes scintillation (fading) and beam wander that require adaptive
optics or diversity. In space there is no weather, and laser inter-satellite links, now flying in large
numbers in low-Earth-orbit constellations (Chapter~\ref{ch:satellite}), carry tens to hundreds of gigabits
per second over thousands of kilometres with telescopes of a few centimetres. Optical feeder links between
satellites and ground stations, and deep-space optical communication, extend the same idea.

%======================================================================
\section{Summary}
\label{sec:ch24:summary}
%======================================================================

\begin{itemize}
\item A transmission line is described by its propagation constant $\gamma=\sqrt{(R+j\omega L)(G+j\omega C)}$
and characteristic impedance $Z_0$. Copper loss in decibels grows as $\sqrt f$ (skin effect) plus a term in
$f$ (dielectric loss); reflections ($\Gamma=(Z_L-Z_0)/(Z_L+Z_0)$) produce echoes and bridged-tap notches;
NEXT grows as $f^{1.5}$ and FEXT as $f^2\ell|H|^2$.
\item The telephone network evolved from manual and electromechanical switching to stored program control
(1965) and digital switching; the hybrid's imperfect 2-wire/4-wire conversion causes echo, which echo
cancellers remove; SS7 moved signalling out of band and made the network programmable.
\item Voiceband modems reached the capacity of the analog voice channel with V.34; V.90 reached 56\,kb/s
by changing the channel model to a digital PCM channel.
\item PDH multiplexed asynchronous tributaries by bit stuffing; SONET/SDH introduced byte-synchronous
125\,$\mu$s frames with pointers, rich overhead and 50\,ms ring protection; OTN wraps any client in frames
with strong FEC.
\item DSL's rate falls steeply with loop length. ADSL used DMT with bit loading over the voice band; VDSL2
from cabinets extended the band to 17--35\,MHz and became FEXT-limited; vectoring cancels FEXT as a
per-tone MIMO problem (precoding downstream, joint detection upstream); G.fast uses TDD to 106--212\,MHz on
loops under 250\,m.
\item HFC cable networks evolved from bonded single-carrier QAM to OFDM with LDPC and 4096-QAM (DOCSIS~3.1),
and to full-duplex echo-cancelled operation or 1.8\,GHz spectrum (DOCSIS~4.0).
\item Twisted-pair Ethernet reached 1 and 10\,Gb/s with four pairs in both directions, echo and NEXT
cancellation, multilevel PAM, trellis or LDPC codes and Tomlinson--Harashima precoding.
\item Single-mode fibre loses about 0.2\,dB/km at 1550\,nm; chromatic dispersion
($D\approx17$\,ps/(nm$\cdot$km)) spreads pulses in proportion to $R_s^2L$ in symbols and causes power
fading in direct detection; PMD grows as $\sqrt L$; the Kerr nonlinearity sets an optimum launch power and
a nonlinear Shannon limit.
\item Optical amplifiers (EDFA, Raman) add ASE with a noise figure of at least 3\,dB; the OSNR budget
$P_{\mathrm{ch}}-L_s-\mathrm{NF}-10\log N+58$\,dB sets the reach of each modulation format. WDM on the
G.694.1 grid and ROADM meshes multiply capacity and flexibility.
\item Coherent receivers detect the optical field, so DSP can compensate CD, PMD and polarisation rotation,
recover carrier phase and decode soft-decision FEC. Symbol rates, constellation shaping and SDM have carried
capacity per wavelength from 100\,Gb/s to over 1\,Tb/s.
\item PONs share a fibre among 32--128 homes with a broadcast downstream and ranged TDMA upstream; their
design is a power budget dominated by the splitter. Data-centre links use PAM-4 IM/DD in the O band, moving
towards linear-drive and co-packaged optics.
\end{itemize}

\begin{labbox}[Wireline and optical links]
The notebook \lab{moderncomms-labs/labs/lab30\_wireline\_optical.py} accompanies this chapter. It builds on
the new module \lab{commlib/wireline.py}, which implements the RLGC cable model, two-port loop analysis with
bridged taps and loading coils, the Unger crosstalk models, gap-approximation DMT bit loading, fibre loss and
dispersion, the IM/DD power-fading response, OSNR and GN-model budgets, and a minimal coherent DSP chain. You
will compute the loss and impedance of a loop and reproduce Figures~\ref{fig:ch24_line_loss}
and~\ref{fig:ch24_loop_impairments}; sweep loop length and noise models to draw your own rate-versus-reach
curves (Figure~\ref{fig:ch24_dsl_rate_reach}); build a small vectored binder with a random FEXT matrix and
compare linear precoding with no cancellation; propagate an IM/DD PAM-4 signal and a coherent DP-QPSK signal
through dispersive fibre and see why one fails and the other does not; implement the CMA butterfly and
Viterbi--Viterbi carrier recovery to reproduce Figure~\ref{fig:ch24_coherent_dsp}; and plan an amplified link
by finding its optimum launch power and margin. The figure script \lab{book/figscripts/ch24\_figs.py}
generates every data figure in this chapter.
\end{labbox}

\problems
\begin{enumerate}[label=\thechapter.\arabic*]
\item A 24 AWG pair has, at 2\,MHz, $R=520$\,$\Omega$/km, $L=0.50$\,mH/km, $C=50$\,nF/km and negligible $G$.
Compute $Z_0$, the attenuation in dB/km and the phase velocity, first with the high-frequency approximations
of~\eqref{eq:ch24:hfloss} and then exactly. How long a loop gives 60\,dB of loss at 2\,MHz?
\item Show from~\eqref{eq:ch24:gammaz0} that in the low-frequency limit $\omega L\ll R$, $G=0$, the
attenuation in dB is proportional to $\sqrt f$ and the phase of $Z_0$ is $-45^\circ$. Use the result to
explain why the step response of a long cable rises in a time proportional to $RC\ell^2$.
\item An open-circuited bridged tap of length $\ell_t$ is connected midway along a loop. Using the input
admittance $\tanh(\gamma\ell_t)/Z_0$, show that the notches occur at $f_k=(2k+1)v/(4\ell_t)$ when loss is
neglected. For $v=0.65c$, where are the first two notches of a 60\,m tap? Which DSL systems would they affect?
\item A hybrid has a balance network of 900\,$\Omega$ in series with 2.16\,$\mu$F. The line presents
$600-j300$\,$\Omega$ at 1\,kHz. What is the transhybrid loss? If the far end of a circuit with 30\,ms one-way
delay returns the echo with this loss, plus 6\,dB of circuit loss in each direction, is an echo canceller
needed according to the 25\,ms rule of thumb?
\item Using the NEXT model~\eqref{eq:ch24:next}, find the NEXT loss at 80\,kHz and at 1\,MHz for 10 and 49
disturbers. Explain why ISDN could use echo-cancelled full duplex across its whole band, but ADSL could not
do so above about 138\,kHz.
\item Compute the line rate of an STS-1, the payload capacity after transport and path overhead, and the
fraction of an OC-48 that carries user data when it transports 48 DS3s. How many 64\,kb/s channels does an
OC-192 carry when it is filled with DS1s in VT1.5s?
\item Explain positive justification in a DS2 multiplexer. If a DS1 tributary runs at $+50$\,ppm and the
multiplexer reserves one stuffing opportunity per tributary in each 1176-bit DS2 frame, what fraction of the
stuffing opportunities carry data? Repeat for $-50$\,ppm.
\item Using the gap approximation with $\Gamma=12$\,dB, a transmit PSD of $-40$\,dBm/Hz, $-140$\,dBm/Hz noise
and a loop loss of $25.4\sqrt{f/1\,\text{MHz}}$\,dB per km, estimate the highest usable ADSL2+ tone and the
approximate downstream rate on a 2.5\,km loop. Compare with Figure~\ref{fig:ch24_dsl_rate_reach}.
\item For a two-line binder with tone-$k$ channel $\vect H=\begin{bmatrix}h_{11}&h_{12}\\h_{21}&h_{22}\end{bmatrix}$,
write the zero-forcing precoder $\vect P$ such that $\vect H\vect P$ is diagonal and equal to
$\mathrm{diag}(h_{11},h_{22})$. If $|h_{12}/h_{11}|=|h_{21}/h_{22}|=0.1$, by how much does precoding raise
the transmit power? What happens when the ratios approach 1, as in G.fast at 200\,MHz?
\item A DOCSIS~3.1 downstream OFDM channel of 96\,MHz uses 25\,kHz subcarriers and a 2.5\,$\mu$s cyclic prefix
with 2048-QAM on all data subcarriers. Estimate its payload rate, assuming 4\% of subcarriers are pilots or
guard and an LDPC code of rate 0.89.
\item Single-mode fibre at 1550\,nm has $D=17$\,ps/(nm$\cdot$km). For IM/DD, find the distance at which the
first power-fading null~\eqref{eq:ch24:fading} falls at half the symbol rate of (a) 10\,Gb/s NRZ, (b)
28\,GBd PAM-4, (c) 56\,GBd PAM-4. What changes at 1310\,nm with $D=1$\,ps/(nm$\cdot$km)?
\item Derive~\eqref{eq:ch24:osnr} from the ASE power $\mathrm{NF}\cdot h\nu G B_{\mathrm{ref}}$ per amplifier.
A 3000\,km link has 100\,km spans of 0.2\,dB/km fibre with 3\,dB of extra loss per span and 5.5\,dB NF
amplifiers. What launch power per channel gives an OSNR of 15\,dB? Is that power likely to be in the
nonlinear regime?
\item Starting from~\eqref{eq:ch24:gn}, derive the optimum launch power and show that at the optimum the
nonlinear interference equals half the ASE. Show that the optimum SNR falls by 2\,dB for every doubling of
the number of spans, and by how much it improves if the amplifier noise figure is reduced by 1\,dB.
\item A GPON uses a 1:64 split with 3\,dB excess loss, 15\,km of fibre at 0.35\,dB/km (upstream), 6
connectors at 0.4\,dB and a 2\,dB repair margin. Which budget class is needed? What is the maximum reach of
the same design with a 1:32 split and class B+ optics?
\item \emph{(Simulation)} Using \lab{commlib/wireline.py}, simulate 32\,GBd DP-QPSK over 1000\,km of fibre
with CD, a random polarisation rotation and laser phase noise. Find the number of CD-compensation taps
needed at 2 samples per symbol, compare CMA and decision-directed LMS convergence in the butterfly, and
measure the BER against SNR with Viterbi--Viterbi carrier recovery for linewidths of 100\,kHz and 1\,MHz.
\item \emph{(Simulation)} Build a 16-line vectored VDSL2 binder: draw a random FEXT matrix per tone whose
magnitudes follow~\eqref{eq:ch24:fext} with random phases, compute each line's rate without vectoring, with
full linear precoding, and with partial vectoring that cancels only the four strongest crosstalkers per tone.
Plot the rate against the number of cancelled crosstalkers.
\end{enumerate}

\furtherreading
\begin{itemize}
\item T.~Starr, J.~M.~Cioffi and P.~J.~Silverman, \emph{Understanding Digital Subscriber Line Technology},
Prentice Hall, 1999. The standard introduction to the loop plant, cable models, crosstalk and DMT; see also
the follow-up by Starr, Sorbara, Cioffi and Silverman, \emph{DSL Advances} (2003).
\item G.~Ginis and J.~M.~Cioffi, ``Vectored transmission for digital subscriber line systems,'' \emph{IEEE
Journal on Selected Areas in Communications}, vol.~20, no.~5, 2002. The paper behind vectoring.
\item J.~Bellamy, \emph{Digital Telephony}, 3rd ed., Wiley, 2000. The telephone plant, switching,
signalling, PDH and SONET from an engineer who knew the Bell System from the inside.
\item G.~P.~Agrawal, \emph{Fiber-Optic Communication Systems}, 4th ed., Wiley, 2010 (5th ed., 2021). The
standard textbook on fibre, sources, detectors, amplifiers, dispersion and system design; his
\emph{Nonlinear Fiber Optics} covers the Kerr effects in depth.
\item K.~C.~Kao and G.~A.~Hockham, ``Dielectric-fibre surface waveguides for optical frequencies,''
\emph{Proceedings of the IEE}, vol.~113, no.~7, 1966. The paper that started optical communication.
\item S.~J.~Savory, ``Digital coherent optical receivers: algorithms and subsystems,'' \emph{IEEE Journal of
Selected Topics in Quantum Electronics}, vol.~16, no.~5, 2010; and K.~Kikuchi, ``Fundamentals of coherent
optical fiber communications,'' \emph{Journal of Lightwave Technology}, vol.~34, no.~1, 2016. Two clear
tutorials on the coherent receiver and its DSP.
\item R.-J.~Essiambre, G.~Kramer, P.~J.~Winzer, G.~J.~Foschini and B.~Goebel, ``Capacity limits of optical
fiber networks,'' \emph{Journal of Lightwave Technology}, vol.~28, no.~4, 2010. The definitive study of the
nonlinear Shannon limit.
\item P.~Poggiolini, ``The GN model of non-linear propagation in uncompensated coherent optical systems,''
\emph{Journal of Lightwave Technology}, vol.~30, no.~24, 2012. The model behind~\eqref{eq:ch24:gn} and
Figure~\ref{fig:ch24_osnr_nli}.
\item P.~J.~Winzer, D.~T.~Neilson and A.~R.~Chraplyvy, ``Fiber-optic transmission and networking: the
previous 20 and the next 20 years,'' \emph{Optics Express}, vol.~26, no.~18, 2018. A superb historical and
forward-looking survey, including the case for SDM.
\item Standards: ITU-T G.992.5 (ADSL2+), G.993.2 and G.993.5 (VDSL2 and vectoring), G.9701 (G.fast), G.707
(SDH), G.709 (OTN), G.652 (single-mode fibre), G.694.1 (DWDM grid), G.984 and G.9807.1 (GPON, XGS-PON);
CableLabs CM-SP-PHYv3.1 and CM-SP-PHYv4.0 (DOCSIS PHY); IEEE 802.3 (Ethernet); OIF 400ZR implementation
agreement. The primary sources for every parameter quoted in this chapter.
\end{itemize}
